% concatenated from subset.tex, paper_macros.tex
%\documentclass[11pt]{article}
\documentclass[US-letter,USenglish,authorcolumns,autoref]{lipics-v2021}
%PACKAGES
%\usepackage[margin=0.95in]{geometry}
\usepackage[noend]{algpseudocode}
\usepackage{algorithm}
\usepackage{times}
\usepackage{amssymb}
\usepackage{amsmath}
\usepackage{latexsym}
\usepackage{graphics}
\usepackage{url}
\usepackage{verbatim}
\usepackage{color}
\usepackage{setspace}
\usepackage{multirow}
%\usepackage{ctable}
\usepackage{lineno}
\usepackage{booktabs}
\usepackage{graphicx}
\usepackage{caption}
\usepackage[shortlabels]{enumitem}

\renewcommand{\tt} {\mathtt}

\definecolor{winered}{rgb}{0.5,0.2,0}
\definecolor{lipicsYellow}{rgb}{0.99,0.78,0.07}

\usepackage{hyperref}
\hypersetup{ 
    colorlinks = true,
    allcolors={winered},
%    linkbordercolor = {white}
}

\usepackage{wrapfig}
\usepackage[normalem]{ulem}
\usepackage{changepage}


%ALGO MACROS
\algnewcommand\And{\textbf{and}}
\algnewcommand\Or{\textbf{or}}
\algnewcommand\Not{\textbf{not}}
\algnewcommand\In{\textbf{in}}
\algnewcommand\Each{\textbf{each}}

\newcommand{\true}{\textsc{true}}
\newcommand{\false}{\textsc{false}}

%THEOREMS
%\newtheorem{theorem}{Theorem}[section]          % Numbers theorems bysection.
%\newtheorem{proposition}[theorem]{Proposition} % Numbers a proposition by a theorem number.
%\newtheorem{lemma}[theorem]{Lemma}             % Numbers a lemma by a theorem number.
%\newtheorem{corollary}[theorem]{Corollary}         % Numbers a corollary by a theorem number.
%\newtheorem{example}[theorem]{Example}         % Numbers an example by a theorem number.
%\newtheorem{examples}[theorem]{Examples}      % Numbers a list of examples by a theorem number.
%\newtheorem{remark}[theorem]{Remark}           % Numbers remark by a theorem number.
%\newtheorem{remarks}[theorem]{Remarks}         % Numbers a list of remarks by a theorem number.
%\newtheorem{conjecture}[theorem]{Conjecture}  % Numbers a conjecture by a theorem number.
%\newtheorem{definition}[theorem]{Definition}      % Numbers a definition by a theorem number.
\newtheorem{fact}[theorem]{Proposition}   			% Numbers a definition by a theorem number.
%\newtheorem{observation}[theorem]{Observation}      % Numbers a definition by a theorem number.
\newtheorem{assumption}[theorem]{Assumption}      % Numbers a definition by a theorem number.
\newtheorem{property}[theorem]{Property} 

%PROOF
%\newenvironment{proof}{{\em Proof.}}{\hspace*{\fill}$\Box$\par\vspace{3mm}}

%SQUISHLISTS
\newcommand{\squishlist}{
 \begin{list}{$\bullet$}
  { \setlength{\itemsep}{0pt}
     \setlength{\parsep}{3pt}
     \setlength{\topsep}{3pt}
     \setlength{\partopsep}{0pt}
     \setlength{\leftmargin}{2.5em}
     \setlength{\labelwidth}{1em}
     \setlength{\labelsep}{0.5em} } }

\newcommand{\squishlisttwo}{
 \begin{list}{$\triangleright$}
  { \setlength{\itemsep}{0pt}
     \setlength{\parsep}{0pt}
    \setlength{\topsep}{0pt}
    \setlength{\partopsep}{0pt}
    \setlength{\leftmargin}{2em}
    \setlength{\labelwidth}{1.5em}
    \setlength{\labelsep}{0.5em} } }

\newcommand{\squishend}{
  \end{list}  }

%LINE SPACING
\newcommand{\ra}[1]{\renewcommand{\arraystretch}{#1}}

\usepackage{listings}

\definecolor{verbgray}{gray}{0.9}

\lstnewenvironment{code}{%
  \lstset{backgroundcolor=\color{verbgray},
 frame=single,
  language=C,
  framerule=0pt,
  basicstyle=\ttfamily,
  showstringspaces=false,
  commentstyle=\color{blue}\textit,
  keywordstyle=\color{black}\bf,
  numbers=none,
  keepspaces=true,
  %numberstyle=\footnotesize, 
  columns=fullflexible}}{}


%EXAMPLE BOXES
%\usepackage[dvipsnames]{xcolor}

\newcommand{\same}[1]{\mathcal{S}_{#1}}
\newcommand{\opposite}[1]{\mathcal{O}_{#1}}
\newcommand{\sala}[1]{\mathcal{D}_{#1}}

\newcommand{\LC}[1]{\mathcal{L}_{#1}}

\newcommand{\RLPCR}{RL-PCR-rep }
\newcommand{\RLPCRs}{RL-PCR-reps }
\newcommand{\RLCCR}{RL-CCR-rep }
\newcommand{\RLCCRs}{RL-CCR-reps }
\newcommand{\RLPRR}{RL-PRR-rep }
\newcommand{\RLPRRs}{RL-PRR-reps }


\newcommand{\PCR}{\textrm{PCR}} 
\newcommand{\CCR}{\textrm{CCR}}
\newcommand{\PRR}{\textrm{PRR}}

\newcommand{\PCRrep}[1]{\mathbf{{RLPCR}}({#1})}
\newcommand{\CCRrep}[1]{\mathbf{{RLCCR}}({#1})}

\newcommand{\RLrep}[1]{\mathbf{{RL}}({#1})}
\newcommand{\LCrep}[1]{\mathbf{{LC}}({#1})}
\newcommand{\Samerep}[1]{\mathbf{{Same}}({#1})}

\newcommand{\RLLrep}[1]{\mathbf{{RL2}}({#1})}
\newcommand{\LCCrep}[1]{\mathbf{{LC2}}({#1})}
\newcommand{\OPPrep}[1]{\mathbf{{OPP}}({#1})}


\newcommand{\Pset}{\mathbb{P}}
\newcommand{\Cset}{\mathbb{C}}
\newcommand{\Rset}{\mathbb{R}}

\newcommand{\Special}{\mathbf{SP}}
\newcommand{\SpecialO}{\mathbf{SP2}}

%===========================
% Bold math fonts for sets
%===========================
\newcommand{\bU}{\mathbf{U}}
\newcommand{\bT}{\mathbf{T}}
\newcommand{\bS}{\mathbf{S}}
\newcommand{\bA}{\mathbf{A}}
\newcommand{\bB}{\mathbf{B}}
\newcommand{\sC}{\mathcal{C}}

\newcommand{\bP}{\mathbf{P}}
\newcommand{\bC}{\mathbf{C}}
\newcommand{\bM}{\mathbf{M}}
\newcommand{\bR}{\mathbf{R}}


\newcommand{\bN}{\mathbf{N}}
\newcommand{\BL}{\mathbf{L}}

\newcommand{\bx}{\mathbf{x}}
\newcommand{\by}{\mathbf{y}}
\newcommand{\bz}{\mathbf{z}}
\newcommand{\bp}{\mathbf{p}}

\newcommand{\bs}{\mathbf{s}}
\newcommand{\ba}{\mathbf{a}}
\newcommand{\bb}{\mathbf{b}}
\newcommand{\br}{\mathbf{r}}
\newcommand{\bq}{\mathbf{q}}
\newcommand{\bt}{\mathbf{t}}

%===========================
% COLORS
%===========================

% http://latexcolor.com/

\definecolor{shadecolor}{rgb}{.91, .91, .91}
\definecolor{bordercolor}{rgb}{.8, .8, .6}

%blues
\definecolor{ultramarine}{rgb}{0, 0.125, 0.376}

 \definecolor{arsenic}{rgb}{0.23, 0.27, 0.29}
 \definecolor{beige}{rgb}{0.96, 0.96, 0.86}

%yellows
\definecolor{amber}{rgb}{1.0, 0.75, 0.0}
\definecolor{orange}{rgb}{1.0, 0.49, 0.0}
\definecolor{dandelion}{rgb}{0.94, 0.88, 0.19}

 
 %greens
  \definecolor{indiagreen}{rgb}{0.07, 0.53, 0.03}
  \definecolor{huntergreen}{rgb}{0.21, 0.37, 0.23}


\newcommand{\blue}[1] {\textcolor{blue}{#1}}
\newcommand{\red}[1] {\textcolor{red}{#1}}
\newcommand{\green}[1] {\textcolor{green}{#1}}

\newcommand{\bblue}[1] {{\bf \textcolor{blue}{#1}}}
\newcommand{\rred}[1] {{\bf \textcolor{red}{#1}}}
\newcommand{\ggreen}[1] {{\bf \textcolor{green}{#1}}}

\newcommand{\defo}[1] {\emph{\textcolor{blue}{#1}}}


%===========================


\definecolor{shadecolor}{rgb}{.9, .9, .9}

 \usepackage{framed}
 
 
 %\colorlet{framecolor}{ultramarine}
 % \colorlet{exframecolor}{orange}

 \setlength\FrameRule{1.5pt}
       
\newenvironment{frshaded}{%
    \def\FrameCommand{\fboxrule=\FrameRule\fboxsep=\FrameSep \fcolorbox{lipicsYellow}{shadecolor}}%
    \MakeFramed {\FrameRestore}}%
    {\endMakeFramed}

    \newenvironment{frshaded*}{%
    \def\FrameCommand{\fboxrule=\FrameRule\fboxsep=\FrameSep \fcolorbox{lipicsYellow}{shadecolor}}%
    \MakeFramed {\advance\hsize-\width \FrameRestore}}%
    {\endMakeFramed}
    
    \newcounter{examplecounter}
\newenvironment{exam}{
 \begin{frshaded*}
    \refstepcounter{examplecounter}%
    \noindent
  \textbf{Example \arabic{examplecounter}}%
  \quad
}{%
\end{frshaded*}
}

%----------------

\newenvironment{frshaded2}{%
    \def\FrameCommand{\fboxrule=\FrameRule\fboxsep=\FrameSep \fcolorbox{lipicsLineGray}{shadecolor}}%
    \MakeFramed {\FrameRestore}}%
{\endMakeFramed}

\newenvironment{frshaded2*}{%
    \def\FrameCommand{\fboxrule=\FrameRule\fboxsep=\FrameSep \fcolorbox{lipicsLineGray}{shadecolor}}%
    \MakeFramed {\advance\hsize-\width \FrameRestore}}%
    {\endMakeFramed}



\newenvironment{result}{
 \begin{frshaded2*}
}{%
\end{frshaded2*}

}

%----------------

\newenvironment{frshaded3}{%
    \def\FrameCommand{\fboxrule=\FrameRule\fboxsep=\FrameSep \fcolorbox{lipicsYellow}{white}}%
    \MakeFramed {\FrameRestore}}%
{\endMakeFramed}

\newenvironment{frshaded3*}{%
    \def\FrameCommand{\fboxrule=\FrameRule\fboxsep=\FrameSep \fcolorbox{lipicsYellow}{white}}%
    \MakeFramed {\advance\hsize-\width \FrameRestore}}%
    {\endMakeFramed}


\newenvironment{result2}{
 \begin{frshaded3*}
}{%
\end{frshaded3*}

}

\graphicspath{{graphics/}}

\DeclareMathOperator{\wt}{weight}
\DeclareMathOperator{\ap}{ap}
\DeclareMathOperator{\neck}{neck}
\DeclareMathOperator{\rank}{rank}
\DeclareMathOperator{\reflect}{reflect}
\DeclareMathOperator{\comp}{comp}
%\DeclareMathOperator{\max}{max}


\nolinenumbers

\renewcommand{\baselinestretch}{1.05}

%==================================================================
%\title{Decoding universal cycles for the $k$-subsets of an $n$-set}



\title{Decoding universal cycles for $t$-subsets and $t$-multisets by decoding bounded-weight de Bruijn sequences}
\titlerunning{~~Decoding universal cycles}
%==================================================================

\author{Daniel Gabri\'{c}}{University of Guelph, Canada}{}{}{}
\author{Wazed Imam}{University of Guelph, Canada}{}{}{}
\author{Lukas Janik-Jones}{University of Guelph, Canada}{}{}{}
\author{Joe Sawada}{University of Guelph, Canada}{}{}{}

\author{~}{~}{}{}{}
\authorrunning{~}

\Copyright{~}
%\ccsdesc[500]{Mathematics of computing~Discrete mathematics~Combinatorics~Combinatorial algorithms}
%\keywords{orientable sequence, de Bruijn sequence, concatenation tree, cycle-joining, universal cycle}


%==================================================================
\begin{document}

%-----------  TITLE STUFF ---------------

%\title{ Decoding universal cycles for the $k$-subsets of an $n$-set}
%\author{ D. Gabric, W. Imam, L. Janik-Jones, J. Sawada \\
%}
\maketitle

\frenchspacing

\begin{abstract}  \normalsize
A universal cycle for a set $\mathbf{S}$ of combinatorial objects is a cyclic sequence of length $|\mathbf{S}|$ that contains a \emph{representative} of each element in $\mathbf{S}$ exactly once as a substring.
Despite the many universal cycle constructions known in the literature for various sets including $k$-ary strings of length $n$, permutations of order $n$, $t$-subsets of an $n$-set, and $t$-multisets of an $n$-set, remarkably few have efficient decoding (ranking/unranking) algorithms.  In this paper we develop the first polynomial time/space decoding algorithms for bounded-weight de Bruijn sequences for strings of length $n$ over an alphabet of size $k$.  %The ranking \red{(do we want to switch back and forth between ranking and encoding, or do we want to stick with one?)} algorithms require $O(n^3k^2)$ time and the unranking algorithms require $\mathcal{O}(n^4k^2 \log k)$ time, both using polynomial space.  
The results are then applied to decode universal cycles for $t$-subsets and $t$-multisets. 
\end{abstract}



%=====================================================================
%=====================================================================
%=====================================================================
\section{Introduction}  







A \defo{universal cycle} for a set $\mathbf{S}$ of combinatorial objects is a cyclic sequence of length $|\mathbf{S}|$ that contains a \emph{representation} of each element in $\mathbf{S}$ exactly once as a substring.  Unless otherwise stated, assume that the representative of each element $\bs \in \bS$ is simply $\bs$ itself.  In cases where $\bs$ is not a string, then its representative will be clearly articulated.  The following ``fundamental questions'' can be asked of any set $\bS$~\cite{chung}:
%
\smallskip

\begin{adjustwidth}{0.5cm}{0.3cm}
\begin{enumerate}
    \item[(i)] Do universal cycles \emph{exist} for $\bS$?
    \item[(ii)] How many universal cycles for $\bS$ are there?
    \item[(iii)] Can a universal cycle for $\bS$ be  (efficiently) constructed?
    \item[(iv)] Does a universal cycle for $\bS$ exist with properties that make it efficiently \emph{decodable}? In other words, are there efficient ranking and ranking algorithms for a specific universal cycle of $\bS$?
\end{enumerate}
\end{adjustwidth}
%

\smallskip
\noindent
There are numerous results pertaining to the first three questions for a wide
variety of interesting sets $\bS$; see for instance~\cite{fred-nfsr,chung,2009research,godbole2011,etzion-book,dbseq}. The decoding question, though, has proven to be more elusive.  

Let $\bS_k(n)$ denote the set of all length-$n$ strings over $\{\tt{1,2,\ldots, k}\}$.  A universal cycle for $\bS_k(n)$ is known as a \defo{de Bruijn sequence}.  The \defo{weight} of a string is the sum of its symbols. Let $\bS_k(n,w^\uparrow)$ denote the subset of $\bS_k(n)$ containing the strings with weight \emph{at least} $w$.  Let $\bS_k(n,w_\downarrow)$ denote the subset of $\bS_k(n)$ containing the strings with weight \emph{at most} $w$.  Universal cycles for $\bS_k(n,w^\uparrow)$ and $\bS_k(n,w_\downarrow)$ are known as a \defo{bounded-weight de Bruijn sequences}. 
In this paper we present efficient ranking and unranking algorithms for a bounded-weight de Bruin sequence.  The algorithms are applied to efficiently decode universal cycles for $t$-subsets and $t$-multisets.  They are the first known efficient algorithms to decode specific universal cycles for these objects.

\medskip
%================================
%================================
\noindent{\bf Decoding universal cycles.}~~
%
Let $\mathcal{U}=u_1u_2\cdots u_m$ be a universal cycle for a set $\mathbf{S}$ of length-$n$ strings.
The \defo{rank} of a string $\bs \in \bS$ is the starting index where $\bs$ is found in $\mathcal{U}$.  If $\bs$ is a prefix of $\mathcal{U}$ then it has rank 1; if $\bs$ starts at the last symbol of $\mathcal{U}$ and wraps around, then it has rank $|\mathbf{S}|$. Given a rank $r$, the \defo{unranking} process determines the string $\bs \in \bS$ starting at index $r$ in $\mathcal{U}$.  Together, ranking and unranking algorithms allow for a universal cycle to be \emph{decoded}.   Any universal cycle can be decoded by traversing the sequence until a given string, or rank, is discovered; however, this requires $O(|{\bf S}|)$ time.  Similarly, a look-up table can be applied to store the position of a given string, but this requires $\Theta(|{\bf S}|)$  space.  We say a decoding algorithm is \defo{efficient} if the ranking and unranking algorithms run in polynomial time and require polynomial space, with respect to the size of the alphabet and the length of each string $\bs$. %\red{do we want to define efficient in terms of the size of the strings, or in terms of the size of the set $\bS$, e.g., a decoding algorithm is efficient if the ranking and unranking can be done in $o(|\bS|)$ time~\cite{Wilf:1982} Then we could expand on what we want, which is even stricter. }
%
The efficient decoding of a universal cycle is necessary when it is applied to robotic position sensing (vision) applications~\cite{BM,magic}.

 There are numerous de Bruijn sequence constructions that apply a variety of different algorithmic approaches; see, for instance~\cite{fred-nfsr,etzion-book,dbseq}.
However, somewhat surprisingly, only the lexicographically smallest de Bruijn sequence for $\bS_k(n)$ is known to be efficiently decodable for all $n$ and~$k$~\cite{kociumaka,ranking}. We review the decoding approach for this de Bruijn sequence in Section~\ref{sec:granddaddy}.
%
%By applying Lempel's $D$-morphism~\cite{lempel}, 
A recursively constructed de Bruijn sequence can be efficiently decoded in the binary case, but not for all $n$ and $k$~\cite{mitchell-decode,tuliani-decode}. See also~\cite[p.317, ex.97-99]{knuth}.
%
The only other known efficient decoding algorithms are for universal cycles of permutations using a shorthand representation~\cite{shorthand,shorthand2}.  This highlights the difficulty in discovering such decoding algorithms.

\medskip
%================================
%================================
\noindent{\bf $t$-subsets and $t$-multisets.}~~
%
Let $\mathbf{Subset}(n,t)$ denote the set of all $t$-subsets of $\{\tt{1,2,\ldots, n}\}$.
Universal cycles for $\mathbf{Subset}(n,t)$  have been considered using a standard string representation for each subset, where, for instance, either $\tt{12}$ or $\tt{21}$ can be used to represent the subset $\{1,2\}$~\cite{chung,jackson,hurlbert,2009research,rudoy}.  Using this representation, universal cycles for $t$-subsets exist only if $n$ divides ${n \choose t}$.  A $t$-subset can also be represented by a length-$n$ binary string with $t$ ones. A shorthand representation excludes the final redundant bit. Applying this shorthand representation, universal cycles for $t$-subsets can be efficiently constructed for all $1 \leq t \leq n$~\cite{fixedweight2}. No efficient decoding algorithms are known for these universal cycles.  
In~\cite{cantwell}, a labelled graph is applied to represent a $t$-subset. With this representation they demonstrate the existence of a universal cycle for $t$-subsets, however, no efficient construction is provided.
A difference representation was recently described in~\cite{difference}.  Given a subset
$\{\tt{s}_1, \tt{s}_2, \ldots , \tt{s}_t\}$ with elements listed in increasing order, its \defo{difference representative} is the length-$t$ string where the first symbol is $\tt{s}_1$ and the $j$-th symbol is  $\tt{s}_j - \tt{s}_{j-1}$ for $1 < j \leq t$. For example, the difference representatives for 
 $\mathbf{Subset}(5,3)$ = \[ \{ \tt{\{1,2,3\}, \{1,2,4\}, \{1,2,5\}, \{1,3,4\}, \{1,3,5\}, \{1,4,5\}, \{2,3,4\}, \{2,3,5\}, \{2,4,5\}, \{3,4,5\}} \}\] 
 are given by
$\tt{111, 112, 113, 121, 122, 131, 211, 212, 221},$ and $\tt{311}$, respectively.  In general, the difference representatives for $\mathbf{Subset}(n,t)$ correspond to the strings in $\bS_{n-t+1}(t,n_\downarrow)$.

%====================
Let $\mathbf{Multiset}(n,t)$ denote the set of all $t$-multisets of $ \{\tt{0,1,\ldots, n-1}\}$.  
Universal cycles for $\mathbf{Multiset}(n,t)$  have been considered using a standard string representation for each subset, where, for instance, either $\tt{112}$, $\tt{121}$, or $\tt{211}$ can be used to represent the multiset $\{1,1,2\}$~\cite{jackson,hurlbert2009}.  Using this representation, universal cycles for $t$-multisets exist only if $n$ divides ${n+t-1 \choose t}$.  Like with $t$-subsets, 
we can also apply a difference representation for $t$-multisets~\cite{difference}.   Given a multiset
$\{\tt{m}_1, \tt{m}_2, \ldots , \tt{m}_t\}$ with elements listed in non-decreasing order, its \defo{difference representative} is the length $t$ string where the first symbol is $\tt{m}_1$ and the $j$-th symbol is  $\tt{m}_j - \tt{m}_{j-1}$ for $1 < j \leq t$.
For example, the difference representatives for 
 $\mathbf{Subset}(3,3)$ = \[ \{ \tt{\{0,0,0\}, \{0,0,1\}, \{0,0,2\}, \{0,1,1\}, \{0,1,2\}, \{0,2,2\}, \{1,1,1\}, \{1,1,2\}, \{1,2,2\}, \{2,2,2\}} \}\] 
 are given by
$\tt{000, 001, 002, 010, 011, 020, 100, 101, 110},$ and $\tt{200}$, respectively.  By replacing each symbol $\tt{x}$ with $\tt{x}+1$ in the difference representatives for $\mathbf{Multiset}(n,t)$, the resulting set of strings is $\bS_{n}(t,(n{+}t{-}1)_\downarrow)$.


%in the binary case for a subset of lengths $n$~\cite{mitchell-decode}.  Subsequently, this method was improved for all values of $n$ in the binary case, and in some non-binary cases~\cite{tuliani-decode}.  There is only one de Bruijn sequence that can be decoded for all $k$ and $n$, the lexicographically smallest de Bruijn sequence~\cite{kociumaka,ranking}.  It can be constructed by 
%it can be generated by a simple Prefer-smallest greedy construction and also by concatenating necklace concatenation approach. The study of the latter construction  leads to efficient decoding algorithms.

%(1) Lex smallest DB ~\cite{kociumaka,ranking} (2) Permutations (cool-lex) ~\cite{shorthand2}.  Also, method for inter-leaved or inductive constructions~\cite{mitchell-decode,tuliani-decode}; see also~\cite[p.317, ex.97-99]{knuth}. 

%Until now, there was no known algorithm to decode a universal cycle for $\mathbf{S}_3(5)$ that runs in polynomial time.

%as noted in ~\cite{tuliani-decode}, the algorithm ``is of limited practical value as it is applicable for only a small subset of the possible values of $n$''. 
% which was re-interpreted  by Knuth~\cite[p.317, ex.97-99]{knuth}. An interleaved de Bruijn sequence has been decoded by Tuliani~\cite{tuliani-decode}.  The only other known decodable de Bruijn sequence is one that has many interesting properties: (1) it is the lexicographically smallest (2) it can be generated by a simple Prefer-smallest greedy construction and (3) it can constructed by concatenating the aperiodic-prefixes of necklaces in lexicographic order.  It is the final property that leads to efficient decoding by applying ranking and uranking algorithms for necklaces. The only other objects that have a known efficient decoding algorithm are permutations which are discussed in X and Y.


%For other combinatorial objects, like subsets, permutations, multi-set permutations, partitions graphs, there is also a dearth of efficient decoding algorithms for any respective universal cycles.  In fact, the only other known efficient decoding algorithms are for two universal cycles for permutations, using a shorthand represenation.



\bigskip

%=====================================================================
\noindent
{\bf Main results.} 
The main results of this paper are as follows:

\begin{enumerate}
    \item We demonstrate that the lexicographically smallest universal cycle for $\bS_k(n,w^\uparrow)$ (a bounded-weight de Bruijn sequence) is efficiently decodable. 
    \item We demonstrate a universal cycle for $t$-subsets that is efficiently decodable.
    \item We demonstrate a universal cycle for $t$-multisets that is efficiently decodable.
   % \item We efficiently rank and unrank the list of $k$-ary necklaces of length $n$ with weight at least $w$ as they appear in lexicographic order.
\end{enumerate}

    


%\begin{theorem}
%    The lexicographically smallest bounded-weight de Bruijn sequence for $\bS_k(n,w)$
%    can be ranked in $\mathcal{O}(n^3k^2)$ time and unranked in $\mathcal{O}(n^4k^2 \log k)$ time, using $\mathcal{O}(n^3k)$ space. 
%\end{theorem}

%
\noindent
As we mention later in the paper, a decodable universal cycle for $\bS_k(n,w^\uparrow)$ can be easily applied to decode a universal cycle for the strings in $\bS_k(n,w_\downarrow)$.
The analysis of all algorithms described in this paper assume the unit-cost RAM model of computation. A C implementation of our decoding algorithms is available at \url{https://debruijnsequence.org/db/subset_decode}.





\bigskip

%=====================================================================
\noindent
{\bf Outline.}
The remainder of the paper is organized as follows.  In Section~\ref{sec:back} we present background definitions and notation.  In Section~\ref{sec:granddaddy} we review the efficient decoding approach for the lexicographically smallest de Bruijn sequence. In Section~\ref{sec:bounded} we generalize these results by considering a lower bound on the weight of the strings;  in Section~\ref{sec:upper} we consider an upper bound on the weight.
%
In Section~\ref{sec:subset}, we apply these results to demonstrate efficiently decodable universal cycles for $t$-subsets and $t$-multisets.  In Section~\ref{sec:rank}, we prove the technical results. 

%In Section~\ref{sec:rank}, we give a brief summary and provide details of future research. 

%which also lead to efficient ranking and unranking algorithms for $k$-ary necklaces of length $n$ with weight at least $w$ as they appear in lexicographic order.  


%=====================================================================
%=====================================================================
%=====================================================================
\section{Preliminaries} \label{sec:back}

Recall that $\bS_k(n)$ denotes the set of all length-$n$ strings over $\{\tt{1,2, \ldots, k}\}$.  Consider two strings $\bs$ and $\bt$.  
Let  $\bs \cdot \bt$ denote the  concatenation of $\bs$ and $\bt$, and   let $\bt^j$ denote $j$ copies of $\bt$ concatenated together.
%
%
A \defo{necklace} is the lexicographically smallest string in an equivalence class of strings under rotation.  
%
Let $\neck(\bs)$ denote the lexicographically smallest rotation of $\bs$, i.e., its necklace. Let $\ap(\bs)$ denote the \defo{aperiodic prefix} of $\bs$, i.e., the shortest string $\bt$ such that $\bs = \bt^j$ for some $j \geq 1$.
For example, $\ap(\tt{1111}) = \tt{1}$.  A string $\bs$ is said to be \defo{aperiodic} (primitive) if $\ap(\bs) = \bs$; otherwise $\bs$ is \defo{periodic} (a power). %A \defo{Lyndon word} is an aperiodic necklace. 


Let $\bN_k(n) = (\sigma_1, \sigma_2, \ldots, \sigma_t)$ denote the tuple consisting of all necklaces in $\bS_k(n)$ listed in lexicographic order. For example, 
%
\[ \bN_2(4) = ( \tt{1111, 1112, 1122, 1212, 1222, 2222} ). \]
%
Amazingly, by concatenating together the aperiodic prefixes of the necklaces in $\bN_k(n)$~\cite{fredricksen-beads-colors} we obtain the lexicographically smallest de Bruijn sequence for $\bS_k(n)$~\cite{fredricksen-beads-colors,fred-lex-least}.
Let
\[\mathcal{G}_k(n) = \ap(\sigma_1) \ap(\sigma_2)  \cdots \ap(\sigma_t), \]
denote this lexicographically smallest de Bruijn sequence for $S_k(n)$.
For example,
%
\[ \mathcal{G}_2(4)  = \tt{1 \cdot 1112 \cdot  1122 \cdot 12 \cdot 1222 \cdot 2}.\]
%
The sequence $\mathcal{G}_k(n)$ is also known as the \defo{Granddaddy}\footnote{
This name was originally given to Martin's prefer-largest greedy construction~\cite{martin} by Knuth~\cite[p.317]{knuth}.  
It is equivalent to $\mathcal{G}_k(n)$ by a simple mapping of the alphabet symbols.} de Bruijn sequence and the \defo{Ford} sequence\footnote{Based on a 1957 technical report by Lester Ford who independently discovered the sequence via an equivalent greedy construction.}. 

\medskip
\noindent
{\bf A note on notation.} In this paper we use bold lower case letters such as $\br,\bs, \bt$ to denote arbitrary strings, while we reserve Greek letters such as $\alpha, \beta,\sigma, \gamma$ to denote necklaces.

%

\subsection{Bounding the weight}
The \defo{weight} of a string $\bs$, denoted by $\wt(\bs)$, is the sum of it symbols.  Recall that $\bS_k(n,w^\uparrow)$ denotes the subset of $\bS_k(n)$ containing strings with weight \emph{at least} $w$.
Let $\bN_k(n,w^\uparrow) = (\alpha_1,\alpha_2, \ldots, \alpha_{m})$ denote the tuple consisting of all necklaces in $\bS_k(n,w^\uparrow)$ listed in lexicographic order.  For example, 
%
\[ \bN_2(4,6^\uparrow) =  ( \tt{1122, 1212, 1222, 2222} ). \]
%

\noindent
Remarkably, the Granddaddy construction can be generalized to $\bS_k(n,w^\uparrow)$.
Let
\[\mathcal{G}_k(n,w^\uparrow) = \ap(\alpha_1) \ap(\alpha_2)  \cdots \ap(\alpha_{m}).  \] 
%
The sequence $\mathcal{G}_k(n,w^\uparrow)$ is the lexicographically smallest bounded-weight de Bruijn sequence for $\bS_k(n,w^\uparrow)$~\cite{weight1,generalize-classic-greedy}.
For example,
%
\[ \mathcal{G}_2(4,6^\uparrow)  = \tt{ 1122 \cdot 12 \cdot 1222 \cdot 2}.\]
%

\noindent
Unfortunately, adapting the Granddaddy for strings with an upper bound on the weight does not necessarily lead to a universal cycle.  For example,  the sequence $\tt{1 \cdot 1112 \cdot 1122 \cdot 12}$ is not a universal cycle for $\bS_2(4,6_\downarrow)$: the string $\tt{2211}$ is missing.

There are a number of other known constructions for bounded-weight de Bruin sequences (for instance, see~\cite{weight1,generalize-classic-greedy,karyframework,concattree}); however, no efficient decoding algorithms for them was previously known. Efficient universal cycles are also known for the subsets of $\bS_k(n)$ with both a lower and upper bound on weights~\cite{sawada2013}.  

\subsection{Necklace properties}

The following five properties of necklaces can be derived from the definition of a necklace.

\begin{property} \label{prop:lastnon}
If $\alpha = \tt{a}_1\cdots \tt{a}_{j-1} \tt{x} \tt{k}^{n-j}$ is in $\bN_k(n)$ for some $j\geq 1$ and symbol $\tt{x} \neq \tt{k}$, then $\tt{a}_1\cdots \tt{a}_j (\tt{x}+1) \tt{k}^{n-j}$ is also a necklace.
\end{property}
 
%
\begin{property} \label{prop:first}
Let $n,k > 1$ and $k < w < kn$.  The first necklace in $\bN_k(n,w^\uparrow)$ is $\tt{1}^{n-j-1}\tt{x}\tt{k}^{j}$,
where $j=\lfloor \frac{w-n}{k-1}\rfloor$ and $\tt{x} = w-(n{-}j{-}1)-kj$; it has weight $w$ and is aperiodic.
\end{property}
%
\begin{property} \label{prop:last}
 Let $n,k > 1$ and $w < kn$.  The last two necklaces in  $\bN_k(n,w^\uparrow)$ are
 $(\tt{k{-}1})\tt{k}^{n-1}$ and $\tt{k}^n$.
\end{property}
%

\noindent
The next property follows from the definition of a periodic necklace.

\begin{property} \label{prop:periodic}
Let $\alpha_i$ and $\alpha_{i+1}$ be consecutive necklaces in $\bN_k(n,w^\uparrow)$ where $n,k > 1$ and $w < kn$.   If $\alpha_i$ is periodic then $\ap(\alpha_i) \cdot \ap(\alpha_{i+1})$ has prefix $\alpha_i$.  Moreover $\alpha_{i+1}$ is aperiodic.
\end{property}
%
\begin{proof}
Since $\alpha_i$ is periodic, it can be written as $\alpha_i = (\tt{a}_1\cdots \tt{a}_p)^j$ for some $j > 1$.  Let $\ell$ be the largest index in $[1,p]$ such $\tt{a}_\ell \neq \tt{k}$.  Such an index exists since $\alpha_i$ is not the last necklace $\tt{k}^n$ in $\bN_k(n,w^\uparrow)$. By Property~\ref{prop:lastnon}, there exists a necklace with prefix $(\tt{a}_1\cdots \tt{a}_p)^{j-1}$ that is lexicographically larger than $\alpha_i$ and also in $\bN_k(n,w^\uparrow)$. Thus $\alpha_{i+1}$ must also have prefix  $(\tt{a}_1\cdots \tt{a}_p)^{j-1}$ since $\bN_k(n,w^\uparrow)$ is lexicographically ordered.  Thus $\ap(\alpha_i) \cdot \ap(\alpha_{i+1})$ has prefix $\alpha_i$. The fact that $\alpha_{i+1}$ is aperiodic is proved in~\cite[Lemma 5]{generalize-classic-greedy}.
\end{proof}

\noindent The following property is proved in~\cite{ranking}.

 \begin{property} \label{prop:suffix}
Let $\alpha = \tt{a}_1\tt{a}_2\cdots \tt{a}_n$ be a necklace and let $1 \leq i \leq  j \leq n$.
Then $\tt{a}_i \tt{a}_{i+1}\cdots a_{\tt{j}-1} \tt{x}$ is lexicographically larger than  $\alpha$ if $\tt{x} > \tt{a}_j$.  
\end{property}


%=====================================================================
%=====================================================================
%=====================================================================
\section{Ranking the Granddaddy de Bruijn sequence $\mathcal{G}_k(n)$} \label{sec:granddaddy}

In this section we review how the Granddaddy de Bruijn sequence $\mathcal{G}_k(n)$ can be efficiently decoded based on the groundbreaking results of Kociumaka,  Radoszewski, and  Rytter~\cite{kociumaka}.  Then in Section~\ref{sec:bounded}, we adapt the approach to decode the bounded-weight de Bruijn sequence $\mathcal{G}_k(n,w^\uparrow)$.

%

Let $\bs$ be a string in $\bS_k(n)$. It can be written as $\bs = \bp\bq$ where $\bq$ is the longest suffix of $\bs$ such that $\bq\bp$ is a necklace.  For example, if $\bs = \tt{221221}$
then $\bp = \tt{22}$ and $\bq = \tt{1221}$.
Let $\rank(\bs)$ denote the rank of $\bs$ in $\mathcal{G}_k(n)$. 
%
The last $n$ length-$n$ substrings of $\mathcal{G}_k(n)$ are all strings $\bs = \tt{k}^{n-j}\tt{1}^j$ for $0  \leq  j < n$. For these strings $\rank(\bs) = k^n - (n-j)+1$. 
%
For all other strings, there are two key steps involved in the efficient computation of $\rank(\bs)$. 

\medskip 
\noindent
{\bf Step 1:} Determine consecutive necklaces $\beta_1$, $\beta_2$, and $\beta_3$ in $\bN_k(n)$ such that $\ap(\beta_1)\ap(\beta_2)\ap(\beta_3)$ contains $\bs$ as a substring. 
%
The following remark summarizes results from~\cite{fredricksen-beads-colors}.

%
\begin{remark} \label{rem:granddaddy}
     There exist unique consecutive necklaces $\beta_1$, $\beta_2$, and $\beta_3$ in $\bN_k(n)$ (considered cyclically) such that 
    $\bp$ is a suffix of $\beta_1$ and $\bq$ is a prefix of $\beta_2$, i.e., $\bs$ is a substring of $\beta_1\beta_2$.  Moreover, $\bs$ is a substring of $\ap(\beta_1)\ap(\beta_2)\ap(\beta_3)$.
    %, noting $\beta_3$ may be needed only if $\beta_2$ is periodic.
\end{remark}
%  
%In the above remark, if $\bs$ is a necklace, $\bq = \beta_2 = \bs$.  
The necklaces $\beta_1,\beta_2$, and $\beta_3$ from the above remark can be determined in $\mathcal{O}(n^2)$ time using $\mathcal{O}(n)$ space~\cite{kociumaka}.

\medskip

\noindent
{\bf Step 2:} Compute $T_k(n,\sigma_i) = |\ap(\sigma_1)\ap(\sigma_2)\cdots \ap(\sigma_{i-1})|$ which counts the number of strings in $\bS_k(n)$ whose necklaces are lexicographically smaller than $\sigma_i$.  For example, 
\[ T_2(4,\tt{1122}) = |\{\tt{1111}\}| + |\{\tt{1112, 1121, 1211, 2111}\}| = 5. \]
The value $T_k(n,\sigma_i)$ can be computed in $\mathcal{O}(n^2)$ time  using $\mathcal{O}(n^2)$ space~\cite{kociumaka}.  

%

\medskip

Putting these two steps together,

\begin{center}
$\rank(\bs) = \left\{ \begin{array}{ll}
         k^n-(n{-}j)+1,    \    &\ \  \mbox{if $\bs = \tt{k}^{n-j}\tt{1}^j$ for some $0 \leq j < n$;}\\
         T_k(n, \beta_2), - |\bp|+1, \  &\ \  \mbox{otherwise.} 
         \end{array} \right. $
\end{center}
%
\begin{exam}
Recall $\mathcal{G}_2(4) = \tt{1  111\underline{2   112}2  12  1222  2}$. Consider $\bs = \tt{2112}$. It is a substring in the concatenation of $\beta_1 = \tt{1112}$ and $\beta_2 = \tt{1122}$, where $\bp = \tt{2}$ and $\bq = \tt{112}$.  Since $T_2(4,\tt{1122}) = 5$, the rank of $\bs$ is $5 - 1 + 1 = 5$.  The rank of $\tt{2211}$ is $2^4 - 2 + 1 = 15$.
\end{exam}
%
\noindent
The string at rank $r$ in $\mathcal{G}_k(n)$ can be determined in $\mathcal{O}(n^3 \log k)$ time by applying a binary search and repeatedly applying the ranking algorithm~\cite{kociumaka}.  See also a reinterpretation of these results in~\cite{ranking}.

%=====================================================================
%=====================================================================
%=====================================================================
\section{Decoding the bounded-weight de Bruijn sequence $\mathcal{G}_k(n,w^\uparrow)$} \label{sec:bounded}

In this section we adapt the decoding algorithm for the Granddaddy to decode the bounded-weight de Bruijn sequence $\mathcal{G}_k(n,w^\uparrow)$.   If $w=k$, then $\mathcal{G}_k(n,w^\uparrow) = \mathcal{G}_k(n)$. If $w=kn$, then $\mathcal{G}_k(n,w^\uparrow)$ is the single symbol $\tt{k}$.   
%Recall that $\bN_k(n,w^\uparrow) = (\alpha_1,\alpha_2, \ldots, \alpha_{m})$ is the tuple consisting of all necklaces in $\bS_k(n,w^\uparrow)$ listed in lexicographic order.  
Throughout the remainder of this section assume that $k < w < kn$ and let:
\begin{itemize}
    \item $\alpha_1 = \tt{a}_1\tt{a}_2\cdots \tt{a}_n$ be the first necklace in $\bN_k(n,w^\uparrow)$, 
    \item $\bs = \bp\bq$ be a string in $\bS_k(n,w^\uparrow)$, where $\bq$ is the longest suffix of $\bs$ such that $\bq\bp$ is a necklace, and
    \item  $\beta_1$ and $\beta_2$ be the necklaces defined in Remark~\ref{rem:granddaddy} for $\bs$.
\end{itemize}
%
Recall that $\beta_1$ has suffix $\bp$ and $\beta_2$ has prefix $\bq$.

To adapt {\bf Step 1} from the previous section, observe that necklaces $\beta_1$ and $\beta_2$ may not satisfy the required weight constraint. This is illustrated in the following example.

%------------------------
\begin{exam}  
Consider $n=3$, $k=4$, and $w=9$.  Since 
\[  \bN_4(3,9^\uparrow)  = \tt{144,~ 234,~ 243,~ 244,~ 333,~ 334,~ 344,~ 444}, \]
we have
\[ \mathcal{G}_4(3,9^\uparrow) = 
\tt{14\underline{4\cdot 23}4\cdot
243\cdot
244\cdot
3\cdot
334\cdot
344\cdot
4}.\]
%
Consider $\bs = \tt{423}$ where $\bp = \tt{4}$ and $\bq = \tt{23}$. In $\mathcal{G}_4(3)$, $\bs$ is found as a substring of $\beta_1 \cdot \beta_2 = \tt{22\underline{4 \cdot 23}3}$. However, in $\mathcal{G}_4(3,9^\uparrow)$ it is found as a substring of $14\underline{4 \cdot 23}4$. %Similarly,  $\tt{414}$ is found as a substring of $\tt{134 \cdot 142}$ in $\mathcal{G}_4(3)$, but is found in the wraparound of $\mathcal{G}_4(3,9)$.    
\end{exam}
%------------------------
%

\noindent
To address this challenge, first consider the strings found in the wraparound of $\mathcal{G}_k(n,w^\uparrow)$.  
%
Recall from Property~\ref{prop:first} that $\alpha_1$ is aperiodic. Thus, $\alpha_1$ is a prefix of $\mathcal{G}_k(n,w^\uparrow)$. 
    %it is aperiodic and  the form $1^j\tt{x}k^{n-j-1}$ for some $\tt{x} \in \Sigma_k$ and $j<n$; it is clearly aperiodic. 
    From Property~\ref{prop:last}, $\alpha_{m-1} = (\tt{k{-}1})\tt{k}^{n-1}$ and $\alpha_m = \tt{k}^n$. Thus, $\mathcal{G}_k(n,w^\uparrow)$ has suffix $\tt{k}^n$.  
%
Therefore, if $\bs$ is of the form $\tt{k}^{n-j}\tt{a}_1\cdots \tt{a}_j$ for some $0 \leq j < n$, then it occurs as one of the last $n$ substrings of $\mathcal{G}_k(n,w^\uparrow)$. The rank of such a string $\bs$ is $|\bS_k(n,w^\uparrow)|-(n-j)+1$.  
%
For all other possible strings $\bs$ consider necklaces $\delta_1$, $\delta_2$, and $\delta_3$ such that

%
\begin{itemize}
 \item $\delta_1$ is the \emph{largest} necklace in $\bN_k(n,w^\uparrow)$ that is lexicographically smaller than or equal to $\beta_1$,  
 \item $\delta_2$ is the \emph{smallest} necklace in $\bN_k(n,w^\uparrow)$ that is lexicographically larger than or equal to $\beta_2$, and
 \item $\delta_3$ is the necklace following $\delta_2$ in $\bN_k(n,w^\uparrow)$.
\end{itemize}
%
By their definitions $\delta_1$, $\delta_2$, and $\delta_3$ appear consecutively in $\bN_k(n,w^\uparrow)$.  The necklaces $\delta_2$ and $\delta_3$ are well defined due to the constraints on $\bs$ and $w$.  The fact that $\delta_1$ is well defined follows from the proof of the following lemma. Ultimately, we will show that $\bs = \bp\bq$ is found in $\mathcal{G}_k(n,w^\uparrow)$ as a substring of $\ap(\delta_1)\ap(\delta_2)\ap(\delta_3)$.




%=============================================
\begin{lemma} \label{lem:d1}
    The string $\ap(\delta_1)$ has suffix $\bp$.
\end{lemma}
%
\begin{proof}
First, we show that $\delta_1$ has suffix $\bp$. Recall that necklace $\beta_1$ has suffix $\bp$. Let $j = |\bp|$ and let $\beta_1 = \tt{b}_1\cdots \tt{b}_{n-j}\bp$. If $\bp \neq \tt{k}^j$, then by Property~\ref{prop:lastnon} it must be that $\beta_2$ also has prefix $\tt{b}_1\cdots \tt{b}_{n-j}$ which  is equal to $\bq$ by the definition of $\beta_2$. Thus $\beta_1 = \bq\bp\in\bN_k(n,w^\uparrow)$. Then by the definition of $\delta_1$, we have $\delta_1 = \beta_1$, meaning $\delta_1$ has suffix $\bp$. Now suppose $\bp = \tt{k}^j$.
 If $\beta_1\in \bN_k(n,w^\uparrow)$, then $\delta_1 = \beta_1$ and thus  $\delta_1$ has suffix $\bp$. Otherwise $\delta_1<\beta_1$ and $\beta_1\not\in \bN_k(n,w^\uparrow)$. Suppose, for a contradiction, that $\delta_1$ does not have $\bp$ as a suffix. If the length-$(n-j)$ prefix of $\delta_1$ is $\tt{b}_1\cdots \tt{b}_{n-j}$, then the weight of $\beta_1$ is greater than the weight of $\delta_1$, a contradiction. If the length-$(n-j)$ prefix of $\delta_1$ is not $\tt{b}_1\cdots \tt{b}_{n-j}$, then it must be smaller than $\tt{b}_1\cdots \tt{b}_{n-j}$. But then by Property~\ref{prop:lastnon}, there is a necklace in $\bN_k(n, w^\uparrow)$ that is larger than $\delta_1$ but smaller than $\beta_1$, a contradiction. Thus $\delta_1$ has suffix $\bp$.
%%

Now we show that $\ap(\delta_1)$ has suffix $\bp$. If $\delta_1$ is aperiodic, then we have already proved that $\bp$ is a suffix of $\ap(\delta_1)=\delta_1$. If $\delta_1$ is periodic then it can be written as $\delta_1 = \gamma^j$ for some $j > 1$ where $\ap(\delta_1) = \gamma$.  If $|\bp| \leq |\gamma|$ then $\bp$ is a suffix of $\ap(\delta_1)$. Otherwise, we have that $\ap(\delta_1)\ap(\delta_2)$ has prefix $\delta_1$ by Property~\ref{prop:periodic}. This implies that $\ap(\delta_2)$ begins with $\gamma^{j-1}$. Since $|\gamma| < |\bp|$, we have that $|\bq| = |\delta_1|-|\bp| < |\delta_1| - |\gamma| = |\gamma^{j-1}|$, which means that $\bq$ is a prefix of $|\gamma^{j-1}|$, and is therefore a prefix of $\delta_1$. But then $\delta_1=\bq\bp$ and $\delta_1$ is a cyclic shift of $\bs$, which means that $\bs$ is periodic with the same period as $\delta_1$. This contradicts $\bq$ being the longest suffix of $\bs$ such that $\bq\bp$ is a necklace.  Thus $\bp$ is a suffix of $\ap(\delta_1)$.
\end{proof}
%=============================================


\begin{lemma} \label{lem:d2}
   The string $\ap(\delta_2)\ap(\delta_3)$ has prefix $\bq$.
\end{lemma}
%
\begin{proof}
First, we show that $\delta_2$ has prefix $\bq$. Recall that $\beta_2$ has prefix $\bq$.  Let $j = |\bp|$.  Since $\bs = \bp\bq$ is in $\bS_k(n,w^\uparrow)$ then $\bq\tt{k}^j$ is also in $\bS_k(n,w^\uparrow)$.
Since $\bq\bp$ is a necklace, then  $\bq\tt{k}^j$ is a necklace by repeated application of Property~\ref{prop:lastnon}.  Thus, since $\delta_2$ is the smallest necklace in $\bN_k(n,w^\uparrow)$ that is lexicographically larger than $\beta_2$, $\delta_2$ must have prefix $\bq$.
%%

Now we show that $\ap(\delta_2)\ap(\delta_3)$ has prefix $\bq$. If $\delta_2$ is aperiodic then clearly $\bq$ is a prefix of $\ap(\delta_2)$. 
If $\delta_2$ is periodic, then it cannot be $\tt{k}^n$ by the restriction on $\bs$. Additionally, we have that $\ap(\delta_2)\ap(\delta_3)$ has prefix $\delta_2$ by Property~\ref{prop:periodic}. Thus $\bq$ is a prefix of $\ap(\delta_2)\ap(\delta_3)$.  
\end{proof}

%=============================================

\begin{corollary}
The string $\bs = \bp\bq$ is a substring of $\ap(\delta_1)\ap(\delta_2)\ap(\delta_3)$.
\end{corollary}

%=============================================

To adapt {\bf Step 2} from the previous section, we compute 
$T_k(n,w,\alpha)$ which counts the number of strings in $\bS_k(n,w^\uparrow)$ whose necklaces are lexicographically smaller than a necklace $\alpha$.
Note that $T_k(n,w,\alpha_i) = |\ap(\alpha_1)\ap(\alpha_2)\cdots\ap(\alpha_{i-1})|$. 
%
Let $\rank'(\bs)$ denote the rank of $\bs$ in $\mathcal{G}_k(n,w^\uparrow)$
and let $S_k(n,w^\uparrow) = |\bS_k(n,w^\uparrow)|$.
Putting this all together, we can compute $\rank'(\bs)$ as follows.

\begin{result} \noindent

\vspace{-0.05in}

\begin{center}
$\rank'(\bs) = \left\{ \begin{array}{ll}
         S_k(n,w^\uparrow)-(n-j)+1    \    &\ \  \mbox{if $\bs = \tt{k}^{n-j}\tt{a}_1\cdots \tt{a}_j$ for some $0 \leq j < n$,}\\
         T_k(n, w,\beta_2) - |\bp|+1 \  &\ \  \mbox{otherwise.} 
         \end{array} \right. $
\end{center}

\vspace{-0.1in}
\end{result}
%
\noindent
Note that is not necessary to compute $\delta_2$, since $T_k(n,w,\delta_2) = T_k(n,w,\beta_2)$.



It remains to analyze the running time required to compute $\rank'(\bs)$.  The following recurrence can be applied to enumerate $S_k(n,w^\uparrow)$ for $n,k\geq 0$ and  $0 \leq w \leq kn$:
\begin{equation}
    S_k(n,w^\uparrow) = \sum_{j=1}^{k} S_k(n-1,\max(0,w-j)^\uparrow), 
\end{equation}
with base cases $S_k(n,w^\uparrow) = k^n$ for $w \leq n$ and $S_k(0,w^\uparrow) = 0$ for $w > 0$.
Applying dynamic programming, the value $S_k(n,w^\uparrow)$ can be computed in $\mathcal{O}(n^2k^2)$ time and   $\mathcal{O}(n^2k)$ space.   In Section~\ref{sec:rank}, the following lemma is proved.
%
%-------------------------------
\begin{lemma} \label{lem:T}
Let $n,k > 1$. For any $\alpha \in \bN_k(n)$ the value $T_k(n,w,\alpha)$ can be computed in $\mathcal{O}(n^3k^2)$ time using $\mathcal{O}(n^3k)$ space. 
\end{lemma}
%

\noindent
Recall that $\beta_2$ can be computed in $\mathcal{O}(n^2)$ time using $\mathcal{O}(n)$ space.  Putting this all together we obtain the following theorem.
%If $\wt(\beta_2) \geq w$ then $\delta_2 = \beta_2$.  Otherwise, $\delta_2$ can be obtained by repeatedly incrementing the last non-$\tt{k}$ in $\beta_2$ a total of $w-\wt(\beta_2)$ times:  From Property~\ref{prop:lastnon} the resulting string is a necklace, it clearly has weight $w$, and it is the smallest necklace in $\mathbf{N}_k(n,w)$ that is lexicographically larger than $\beta_2$. By updating the symbol at each index once, this computation can be done in $\mathcal{O}(n)$ time and $\mathcal{O}(n)$ space.  Thus, $\delta_2$ can be computed in $\mathcal{O}(n^2)$ time using $\mathcal{O}(n)$ space.




%-------------------------------
%
%Let $\bs$ be a string in $\bS_k(n,w^\uparrow)$, and let $\bp$ and $\beta_2$ be determined as in Lemma~\ref{lem:alpha2}. 

%
%Note that the case when $\beta_2$ is periodic is handled by Property~\ref{prop:periodic} (like with the Granddaddy) as we illustrate in the following example.  
%---------------
%\begin{exam}
%Recall
%\[ \mathcal{W}_4^{\uparrow}(3,9) = 
%\tt{144\cdot
%234\cdot
%243\cdot
%244\cdot
%3\cdot
%334\cdot
%344\cdot
%4}.\]
%If $\bs = 443$, then $\beta_2 = 333$.  If $\alpha = %433$, then $\beta_2 = 333$.  If $\alpha = 333$, the $\beta_2 = 334$.
%\end{exam}
%---------------

%\noindent
%If $\beta_1$ is periodic, then by the uniqueness of $\beta_1$ and $\beta_2$ together with Property~\ref{prop:periodic}, it can be deduced that $|\bp| \leq |\ap(\beta_1)|$ \red{almost - but I think if $\beta_1 = \tt{x}^n$ there is a problem}.  Thus, $\alpha$ is found in $\mathcal{G}_k(n,w^\uparrow)$ as a substring of $\ap(\beta_1)\ap(\beta_2)\ap(\beta_3)$, where $\beta_3$ follows $\beta_2$ in $\bN_k(n,w^\uparrow)$ (it is needed when $\beta_2$ is periodic). 
%

\noindent


%-------------------------------
\begin{theorem}
The universal cycle $\mathcal{G}_k(n,w^\uparrow)$  can be ranked in $\mathcal{O}(n^3k^2)$ time using $\mathcal{O}(n^3k)$ space.
\end{theorem}
%-------------------------------

\subsection{Unranking}

%Let $1 \leq r \leq |\mathcal{G}_k(n,w^\uparrow)|$.  

In this section, we demonstrate how to determine the string $\bs = \tt{s}_1\tt{s}_2\cdots \tt{s}_n$ at rank $r$ in $\mathcal{G}_k(n,w^\uparrow)$.   We consider two cases depending on the value of $r$.
%

{\bf Case 1:}  $S_k(n,w^\uparrow) - n +1 < r \leq S_k(n,w^\uparrow)$.  In this case, $\bs$ is found in the wraparound.  Recall that $\tt{a}_1\cdots \tt{a}_n$ are the first $n$ symbols in $\mathcal{G}_k(n,w^\uparrow)$ as defined in Property~\ref{prop:first}, and the last $n$ symbols  are $\tt{k}^n$ (from Property~\ref{prop:last}).  Thus, $\bs$ can be returned in $\mathcal{O}(n)$ time.
 %
 
{\bf Case 2 :} $1 \leq r  \leq S_k(n,w^\uparrow) - n +1$.
%
By the construction of $\mathcal{G}_k(n,w^\uparrow)$ and Property~\ref{prop:periodic}, the ranks of the necklaces as they appear in $\bN_k(n,w^\uparrow)$ are strictly increasing.  Let {\sc SmallestNeck}($r$) denote the smallest necklace in $\bN_k(n,w^\uparrow)$ with rank  greater than or equal to $r$ in $\mathcal{G}_k(n,w^\uparrow)$.  If $\gamma_1 =$ {\sc SmallestNeck}($r$), then let $r'$ denote the rank of $\gamma_1$ in $\mathcal{G}_k(n,w^\uparrow)$.  If $r = r'$, then $\bs = \gamma_1$.  Otherwise, let $\gamma_2 = $
{\sc SmallestNeck}($r'-n$). By Property~\ref{prop:periodic}, $\gamma_2$ and $\gamma_1$ appear consecutively in $\bN_k(n,w^\uparrow)$, at most one of the two necklaces is periodic, and $\bs$ is obtained by concatenating the last $r'-r$ elements from $\gamma_2$  with the first $n-(r'-r)$ elements from $\gamma_1$. 
%
\begin{exam}
Let $n=3, k=4$, and $w=9$, and consider $r=3$. Recall 
that  \[ \mathcal{G}_4(3,9^\uparrow) = 
\tt{14\underline{4\cdot 23}4\cdot
243\cdot
244\cdot
3\cdot
334\cdot
344\cdot
4}.\]
 The necklace $\gamma_1 =$ {\sc SmallestNeck}($r$) = $\tt{234}$ has rank $r' = 4$ and
$\gamma_2 =$ {\sc SmallestNeck}$(r'-n) = \tt{144}$. Since $r'-r=1$, the string at rank $r=3$ is the last symbol of $\gamma_2$ concatenated with the first $n-1=2$ symbols of $\gamma_1$. It is  $\bs=423$.


\end{exam}

\noindent  
The necklace {\sc SmallestNeck}($r$) can be computed by making $\mathcal{O}(n \log k)$ rank queries as illustrated in Algorithm~\ref{algo:unrank}.  It follows the same approach as the algorithm in~\cite{ranking} adding the weight constraint.  The algorithm iteratively determines each symbol of $\tau = \bt_1\cdots \bt_n$ starting with the first symbol. The loop invariant after $i$ iterations maintains the property that $\bt_1\cdots \bt_i\tt{k}^{n-i}$ is a necklace with rank greater than or equal to $r$ such that $\bt_1\cdots \bt_{i-1}(\bt_{i}-1)\tt{k^{n-i}}$ is either not in $\bN_k(n,w^\uparrow)$, or it is in $\bN_k(n,w^\uparrow)$ and has rank less than $r$. Binary search, is applied to determine the $i$-th symbol using at most $\log k$ ranking queries.  

%==================================================
\begin{algorithm}[htb]
\small
\caption{Compute the smallest necklace $\bt_1\bt_2\cdots \bt_n$ in $\bN_k(n,w^\uparrow)$ with rank greater than or equal to $r$, where $1 \leq r \leq S_k(n,w^\uparrow) - r + 1$}\label{algo:unrank}
\begin{algorithmic}[1]
\Statex
\Function{SmallestNeck}{$r$}
\Statex

	\State \blue{$\triangleright$ Apply a binary search to find each $\bt_i$}
        \For {$i$ {\bf from} $1$ {\bf to} $n$} 
        		\State $min \gets 1$
		        \State $max \gets k$
                \State $\bt_i \gets \tt{k}$
        		\While{$min < max$}	
			\State $prev \gets \bt_i$
			 \State $v \gets (min + max)/2$
			 \State $\bt_i \gets v$
			 \If{ $\bt_1\cdots \bt_i\tt{k}^{n-i} \in \bN_k(n,w^\uparrow)$ {\bf and} $\rank'(\bt_1\cdots \bt_i\tt{k}^{n-i}) \geq r$}     \ \ $max \gets v$
			 \Else
			 	\State $\bt_i \gets prev$
				\State $min \gets v+1$
			\EndIf
		\EndWhile
	\EndFor	

	\State \Return $\bt_1\bt_2\cdots \bt_n$
	
 \EndFunction
      
\end{algorithmic}
\end{algorithm}



%-------------------------------
\begin{theorem}
The universal cycle $\mathcal{G}_k(n,w^\uparrow)$  can be unranked in $\mathcal{O}(n^4k^2\log k)$ time using $\mathcal{O}(n^3k)$ space.
\end{theorem}
%-------------------------------
\begin{proof}
Testing whether a string is a necklace can be determined in $\mathcal{O}(n)$ time/space~\cite{Booth}.  Thus, the running time and space for the described unranking algorithm is dominated by the $\mathcal{O}(n \log k)$ calls to the ranking function.  
\end{proof}


%=====================================================================
%=====================================================================
%=====================================================================
\section{An upper bound on weight} \label{sec:upper}
Recall that  $\bS_k(n,w_\downarrow)$ denotes the subset of $\bS_k(n)$ containing strings with weight \emph{at most} $w$.
 %In this section, we demonstrate an efficiently decodable universal cycle for $\bS'_k(n,w)$.
Let the \defo{complement} of a string $\bs$, denoted by $\comp(\bs)$, be obtained by interchanging each symbol $\tt{x}$ with $\tt{k{-}x{+}1}$.  
 If $\bs$ has weight $w$, then $\comp(\bs)$ has weight $kn-w+n$.
%
%-------------------------------
\begin{exam}
Let $n=5$, $k=4$, and consider $\bs = \tt{12443}$. It has weight $w=14$.  Then $\comp(\bs) = \tt{43112}$ has weight $20-14+5 = 11$; the symbols $\tt{1}$ and $\tt{4}$ are interchanged, and the symbols $\tt{2}$ and $\tt{3}$ are interchanged.  
\end{exam}
%-------------------------------
%

\noindent
The complement of a universal cycle for $\bS_k(n,(kn-w+n)^\uparrow)$ is a universal cycle for $\bS(n,w_\downarrow)$, and vice versa.  Thus, $\mathcal{G}_k(n,w_\downarrow) = \comp(\mathcal{G}_k(n,kn-w+n)^\uparrow)$ is a universal cycle for $\bS_k(n,w_\downarrow)$.  To rank a string $\bs$ in $\mathcal{G}_k(n,w_\downarrow)$ simply return the rank of $\comp(\bs)$ in $\mathcal{G}_k(n,(kn-w+n)^\uparrow)$.

%-------------------------------
\begin{corollary} \label{cor:S2}
The universal cycle $\mathcal{G}_k(n,w_\downarrow)$  can be ranked in $\mathcal{O}(n^3k^2)$ time and unranked in $\mathcal{O}(n^4k^2\log k)$ time  using $\mathcal{O}(n^3k)$ space.
\end{corollary}
%-------------------------------

\begin{comment}
From~\cite{weight1,concattree}, $\mathcal{G}_k(n,w^\uparrow)$ can be constructed in $\mathcal{O}(1)$-amortized time per symbol using $\mathcal{O}(n^2)$ space.  Thus, the same analysis can be applied to $\mathcal{G}_k(n,w^\uparrow)$ by simply complementing the output of each symbol.

%-------------------------------
\begin{theorem} \label{thm:construct}
Let $n,k > 2$ and $k \leq w \leq nk$.  Then $\mathcal{G}_k(n,w^\uparrow)$ and $\mathcal{W}_k^{\downarrow}(n,w)$ can be constructed in $\mathcal{O}(1)$-amortized time using $\mathcal{O}(n^2)$ space.
\end{theorem}
%-------------------------------
\end{comment}

%=====================================================================
%=====================================================================
%=====================================================================
\section{Applications to $t$-subsets and $t$-multisets} \label{sec:subset}



Recall that $\mathbf{Subset}(n,t)$ denotes the set of all $t$-subsets of $\{\tt{1,2,\ldots , n}\}$. 
Using difference representatives for each subset, a universal cycle for $\mathbf{Subset}(n,t)$ is precisely  a universal cycle for  $\mathbf{S}_{n-t+1}(t,n_\downarrow)$~\cite{difference}.  
Thus, we can apply the results from Corollary~\ref{cor:S2} to efficiently rank a universal cycle for $\mathbf{Subset}(n,t)$.  
%


%-------------------------------
\begin{theorem}
 The sequence $\mathcal{G}_{n-t+1}(t,n_\downarrow)$ is a universal cycle for $\mathbf{Subset}(n,t)$ applying the difference representatives. 
 %It can be generated in $\mathcal{O}(1)$-amortized time per symbol using $O(t^2)$ space.  
 Using $\mathcal{O}(nt^3)$ space, the sequence can be ranked in $\mathcal{O}(n^2t^3)$ time and unranked in $\mathcal{O}(n^2t^4\log n)$ time.
\end{theorem}
%-------------------------------

\begin{exam}
%Consider $\mathbf{Subset}(5,3)$.  
The sequence $\mathcal{G}_3(3,5_\downarrow) = \tt{3112212111}$ is a universal cycle for $\mathbf{Subset}(5,3)$ using difference representatives.  It can be efficiently decoded using its complement
$\mathcal{G}_3(3,7^\uparrow) = \tt{1332232333}$ as illustrated in the following table.
%
\begin{center}
\begin{tabular}{c|c|c|c}
  {\bf Rank}   & {\bf String in $\mathcal{G}_3(3,7^\uparrow)$}  & {\bf Complement (diff rep)} & {\bf Subset over $\{1,2,3,4,5\}$}  \\ \hline
    1   & $\tt{133}$   & $\tt{311}$   & $\{\tt{3,4,5}\}$ \\   
    2   & $\tt{332}$   & $\tt{112}$   & $\{\tt{1,2,4}\}$ \\   
    3   & $\tt{322}$   & $\tt{122}$   & $\{\tt{1,3,5}\}$ \\   
    4   & $\tt{223}$  & $\tt{221}$   & $\{\tt{2,4,5}\}$ \\   
    5   & $\tt{232}$   & $\tt{212}$   & $\{\tt{2,3,5}\}$ \\   
    6   & $\tt{323}$   & $\tt{121}$   & $\{\tt{1,3,4}\}$ \\   
    7   & $\tt{233}$   & $\tt{211}$   & $\{\tt{2,3,4}\}$ \\   
    8   & $\tt{333}$   & $\tt{111}$   & $\{\tt{1,2,3}\}$ \\   
    9   & $\tt{331}$   & $\tt{113}$   & $\{\tt{1,2,5}\}$ \\   
    10  & $\tt{313}$   & $\tt{131}$   & $\{\tt{1,4,5}\}$ \\   
\end{tabular}
\end{center}

\vspace{-0.15in}
 
\end{exam}






Recall that $\mathbf{Multiset}(n,t)$ denotes the set of all $t$-multisets of $\{\tt{0,2,\ldots , n{-}1}\}$. 
Using difference representatives for each multiset and incrementing all symbols by 1, a universal cycle for $\mathbf{Multiset}(n,t)$ is precisely  a universal cycle for  $\mathbf{S}_{n}(t,(n+t-1)_\downarrow)$~\cite{difference}.  
Again, we can apply the results from Corollary~\ref{cor:S2} to efficiently decode a universal cycle for $\mathbf{Multiset}(n,t)$.  
%


%-------------------------------
\begin{theorem}
 The sequence $\mathcal{G}_{n}(t,(n+t-1)_\downarrow)$ is a universal cycle for $\mathbf{Multiset}(n,t)$ applying the adjusted difference representatives. 
 %It can be generated in $\mathcal{O}(1)$-amortized time per symbol using $O(t^2)$ space.  
 Using $\mathcal{O}(nt^3)$ space, the sequence can be ranked in $\mathcal{O}(n^2t^3)$ time and unranked in $\mathcal{O}(n^2t^4\log n)$ time.
\end{theorem}
%-------------------------------




\begin{exam}
%Consider $\mathbf{Subset}(5,3)$.  
The sequence $\tt{2001101000}$ is a universal cycle for $\mathbf{Multieset}(3,3)$ using difference representatives. By incrementing each symbol in the sequence by one, we obtain $\mathcal{G}_3(3,5_\downarrow) = \tt{3112212111}$.
It can be efficiently decoded using its complement
$\mathcal{G}_3(3,7^\uparrow) = \tt{1332232333}$ as illustrated in the following table.
%
\begin{center}
\begin{tabular}{c|c|c|c}
  {\bf Rank}   & {\bf String in $\mathcal{G}_3(3,7^\uparrow)$}  & {\bf Complement minus 1 (diff rep)} & {\bf Multiset over $\{0,1,2\}$}  \\ \hline
    1   & $\tt{133}$   & $\tt{200}$   & $\{\tt{2,2,2}\}$ \\   
    2   & $\tt{332}$   & $\tt{001}$   & $\{\tt{0,0,1}\}$ \\   
    3   & $\tt{322}$   & $\tt{011}$   & $\{\tt{0,1,2}\}$ \\   
    4   & $\tt{223}$  & $\tt{110}$   & $\{\tt{1,2,2}\}$ \\   
    5   & $\tt{232}$   & $\tt{101}$   & $\{\tt{1,1,2}\}$ \\   
    6   & $\tt{323}$   & $\tt{010}$   & $\{\tt{0,1,1}\}$ \\   
    7   & $\tt{233}$   & $\tt{100}$   & $\{\tt{1,1,1}\}$ \\   
    8   & $\tt{333}$   & $\tt{000}$   & $\{\tt{0,0,0}\}$ \\   
    9   & $\tt{331}$   & $\tt{002}$   & $\{\tt{0,0,2}\}$ \\   
    10  & $\tt{313}$   & $\tt{020}$   & $\{\tt{0,2,2}\}$ \\   
\end{tabular}
\end{center}

\vspace{-0.15in}
 
\end{exam}

%=====================================================================
%=====================================================================
%=====================================================================


%=====================================================================
%\section{Ranking necklaces with bounded weight} \label{sec:rank}
\section{Proof of Lemma~\ref{lem:T}: Enumerating $T_k(n,w,\alpha)$} \label{sec:rank}

In this section we demonstrate how to efficiently compute $T_k(n,w,\alpha)$, where $\alpha$ is a necklace in $\bN_k(n,w^\uparrow)$. 
%We also show how this value can be used to efficiently rank and unrank the necklaces in $\bN_k(n,w^\uparrow)$.  
The approach follows the methods as applied in~\cite{ranking,hartman}.
%
To simplify notation going forward assume that $n,k$ and the necklace $\alpha = \tt{a}_1\tt{a}_2\cdots \tt{a}_n$ are fixed. 


%=====================================================================
\subsection{Enumerating preliminary sets} \label{sec:B}

In order to efficiently compute compute $T_k(n,w,\alpha)$, we need to efficiently determine the cardinality of two special sets of strings that have both prefix and suffix restrictions.  For $t \geq j$, let  $\mathbf{B}(t,j,w)$ denote the subset of strings in $\bS_k(t,w)$  where each string has prefix $\tt{a}_1\tt{a}_2\cdots \tt{a}_{j}$ and every non-empty suffix is greater than $\alpha$. 
  When $j=0$, there is no prefix restriction on the strings.

\begin{exam}
Let $\alpha = \tt{112122}$ and $k=3$.  Then the strings in $\mathbf{B}(6,3,11)$ are:

\begin{center}
\begin{tabular}{c c c}
\textcolor{blue}{$\tt{112}$}\textcolor{red}{$\tt{1}$}$\tt{33}$  & \textcolor{blue}{$\tt{112}$}\textcolor{red}{$\tt{2}$}$\underline{\tt{23}}$ & \textcolor{blue}{$\tt{112}$}\textcolor{red}{$\tt{3}$}$\underline{\tt{13}}$ \\
                            & \textcolor{blue}{$\tt{112}$}\textcolor{red}{$\tt{2}$}$\underline{\tt{32}}$ & \textcolor{blue}{$\tt{112}$}\textcolor{red}{$\tt{3}$}$\underline{\tt{22}}$ \\
                            & \textcolor{blue}{$\tt{112}$}\textcolor{red}{$\tt{2}$}$\underline{\tt{33}}$ & \textcolor{blue}{$\tt{112}$}\textcolor{red}{$\tt{3}$}$\underline{\tt{23}}$ \\
                             &  & \textcolor{blue}{$\tt{112}$}\textcolor{red}{$\tt{3}$}$\underline{\tt{32}}$ \\
                             &  & \textcolor{blue}{$\tt{112}$}\textcolor{red}{$\tt{3}$}$\underline{\tt{33}}$ \\
\end{tabular}
\end{center}
%Note that $B(6,3,11) = 9$. 
The string in the first column can be given by $\mathbf{B}(6,4,11)$,
while the  underlined strings in the second and third columns correspond to the sets $\mathbf{B}(2,0,5)$ and $\mathbf{B}(2,0,4)$, respectively.
\end{exam}

\begin{comment}
%-----------------------------
\begin{exam}
Let $\alpha=\T{aaabcc}$ and $\Sigma = \{\T{a,b,c}\}$.   $\mathbf{B}_\alpha(2,0) = \{\T{ab,ac,bb,bc,cb,cc}\} $ consists of strings of length 2 where every non-empty suffix is greater than $\alpha$.  The strings in $\{\T{aa,ba,ca}\}$ are not in this set because the suffix \T{a} is lexicographically smaller than $\alpha$.   Now consider $\mathbf{B}_\alpha(5,2)$ partitioned on the 3rd  (underlined) character: 
%
\vspace{-0.25in}
\begin{center}
\begin{tabular}  {l  @{\hskip 0.4in}  l @{\hskip 0.4in}  l }  \\
\T{aa\blue{\underline{a}}\:cb}  &   \T{aa\blue{\underline{b}}\:ab}   & \T{aa\blue{\underline{c}}\:ab}   \\
\T{aa\blue{\underline{a}}\:cc}  &   \T{aa\blue{\underline{b}}\:ac}   & \T{aa\blue{\underline{c}}\:ac}   \\
                            &   \T{aa\blue{\underline{b}}\:bb}   & \T{aa\blue{\underline{c}}\:bb}   \\
                             &  \T{aa\blue{\underline{b}}\:bc}   & \T{aa\blue{\underline{c}}\:bc}   \\
                             &   \T{aa\blue{\underline{b}}\:cb}   & \T{aa\blue{\underline{c}}\:cb}   \\
                             &   \T{aa\blue{\underline{b}}\:cc}   & \T{aa\blue{\underline{c}}\:cc}.   \\
\end{tabular}
\end{center}

\noindent
Observe how this partition decomposes the set revealing recursive structure:  
\[\mathbf{B}_\alpha(5,2) = \mathbf{B}_\alpha(5,3) \  \cup \  \T{aab} \cdot \mathbf{B}_\alpha(2,0)  \ \cup  \  \T{aac} \cdot \mathbf{B}_\alpha(2,0).\]
%
\end{exam}
%-----------------------------
\end{comment}

%
\noindent
Let $B(t,j,w) = |\mathbf{B}(t,j,w)|$.
An enumeration formula for $B(t,j,w)$, where $0\leq j\leq t \leq n$, can be derived based on the recursive structure illustrated in the previous example.  Considering the empty string, $B(0,0,w) = 1$ if $w \leq 0$ and $B(0,0,w) = 0$ if $w > 0$.  When $t > 0$,  $B(t,t,w) = 0$ because  the suffix $\tt{a}_1\tt{a}_2\cdots \tt{a}_t$ is  lexicographically smaller than or equal to~$\alpha$.  Otherwise, the strings in $\mathbf{B}(t,j,w)$ can be partitioned based on the symbol in position $j{+}1$.
\squishlisttwo
\item If the $j{+}1$st symbol is smaller than $\tt{a}_{j+1}$  then the suffix starting from the first index would be smaller than $\alpha$, a contradiction.  
\item If the $j{+}1$st  symbol is $\tt{a}_{j+1}$ then the number of such strings is $B(t,j+1,w)$.  
\item If the $j{+}1$st  symbol is larger than $\tt{a}_{j+1}$, then any suffix starting at index $1,2,\ldots, j+1$ is 
larger than $\alpha$ by Lemma~\ref{prop:suffix}.  Thus, for each $\tt{x} > \tt{a}_{j+1}$, the remaining $t-j-1$ positions can be filled in $B(t-j-1,0,w-\wt(\tt{a}_1\cdots \tt{a}_t\tt{x}))$ ways.  
\squishend
Thus, for $0\leq j\leq t \leq n$ and $w \geq 0$:
%
\begin{center} \small
$B(t,j,w)  = \left\{ \begin{array}{ll}
    1	 \ & \ \  \   \mbox{if \ \ $j=t=0$ and $w \leq 0$, } \\
    0	 \ & \ \  \  \mbox{if \ \ $j=t$,} \\
    B(t,j+1,w) + \displaystyle{\sum_{\tt{x} = \tt{a}_{j+1}+1}^\tt{k} B(t-j-1,0, w-\wt(\tt{a}_1\cdots \tt{a}_j\tt{x})) } \  &\ \ \  \mbox{otherwise.}
    \end{array} \right.
            \label{eq:B}$
\end{center}
%


%================
Now consider the set $\mathbf{P}(t,j,w)$ which denotes the subset of  $\bS_k(t)$ where each string has weight \emph{exactly} $w$ and prefix  $\tt{a}_1\tt{a}_2\cdots \tt{a}_{j}$ such that every non-empty suffix is greater than $\alpha$. Let $P(t,j,w) = |\mathbf{P}(t,j,w)|$.  Applying the same recursive strategy applied for computing $B(t,j,w)$, but with different base cases,  we obtain the following recurrence for $0\leq j \leq t$:
%
\[
\small 
{P}(t,j,w) = 
\begin{cases}	
	1 & \ \ \  \text{if } j=t=w=0, \\
	0 & \ \ \  \text{if } j=t \mbox{ or } w \leq 0,  \\
	{P}(t,j+1,w) + \displaystyle{
        \sum_{\tt{x}=\tt{a}_{j+1}+1}^{\tt{k}}P(t-j-1,0,w -\wt(\tt{a}_1\cdots \tt{a}_j\tt{x}))} & \ \ \ \text{otherwise}. \\
\end{cases}
\]
\normalsize



%=====================================================================
\subsection{Computing $T_k(n,w,\alpha)$}



Let $\textbf{T}_k(n,w,\alpha)$ denote the subset of strings in $\bS_k(n,w^\uparrow)$ whose necklaces are lexicographically smaller than a given necklace $\alpha$.  Note that $T_k(n,w,\alpha) = |\textbf{T}_k(n,w,\alpha)|$.  To compute $T_k(n,w,\alpha)$, we partition $\textbf{T}_k(n,w,\alpha)$ into subsets following the approaches used in~\cite{ranking,hartman}.  

%
Let $\textbf{A}(t,j)$ denote the subset of strings  in $\textbf{T}_k(n,w,\alpha)$ of the form $\bs= \tt{s}_1\tt{s}_2\cdots \tt{s}_{n}$ where $t$ is the smallest integer such that $\bs'= \tt{s}_t\tt{s}_{t+1}\cdots \tt{s}_1\tt{s}_2\cdots \tt{s}_{t-1} < \alpha$ and $j$ is the largest integer such that $\tt{a}_1 \cdots \tt{a}_j$ is a prefix of $\bs'$.   Let $A(t,j) =  |\textbf{A}(t,j)|$. %
 Then
\[
T_k(n,w, \alpha) =  \displaystyle{\sum_{t=1}^{n}\sum_{j=0}^{n} A(t,j).}
\]

%Consider for some set $\textbf{A}_k(t,j,w, \alpha)$, what would the $j+1$'st element be? Well we know that the string is smaller then or equal to $\alpha$ and that it can't be equal to $\tt{a}_{j+1}$ as that would imply the string is apart $\textbf{A}_k(t,j+1,w, \alpha)$. Thus we can see that the $j+1'st$ element must be smaller than $\tt{a}_{j+1}$.
%EXAMPLE 3

\begin{exam}
\begin{comment}
 The sequence of necklaces $\bN_3(5,10)$ is \small
    \[ 
\tt{1 1 2 3 3,
1 1 3 2 3,
1 1 3 3 2,
1 1 3 3 3,
1 2 1 3 3,
1 2 2 2 3,
1 2 2 3 2,
1 2 2 3 3,
\underline{1 2 3 1 3},
1 2 3 2 3,
1 2 3 3 2,
1 2 3 3 3,
1 3 1 3 2,
1 3 1 3 3, 
1 3 2 2 2,} \] 
%
\vspace{-0.25in}
%
\[
\tt{1 3 2 2 3,
1 3 2 3 2,
1 3 2 3 3,
1 3 3 2 2,
1 3 3 2 3,
1 3 3 3 2,
1 3 3 3 3,
2 2 2 2 2,
2 2 2 2 3,
2 2 2 3 3,
2 2 3 2 3,
2 2 3 3 3,
2 3 2 3 3,
2 3 3 3 3,
3 3 3 3 3. }
\]
\end{comment}
\noindent
The set $\textbf{T}_3(5,10,\tt{12313})$ contains 40 strings from the equivalence classes
represented by the eight necklaces:
$\{ \tt{1 1 2 3 3,
1 1 3 2 3,
1 1 3 3 2,
1 1 3 3 3,
1 2 1 3 3,
1 2 2 2 3,
1 2 2 3 2,
1 2 2 3 3}\}$
 that appear before $\tt{12313}$ in $\bN_3(5,10)$. These strings can be partitioned into 10 subsets of the form $\textbf{A}(t,j)$ 
as follows:

%We can calculate this by calculating the count of each of $\textbf{A}_3(t,j,w,\tt{12313})$ for all $1 \leq t \leq n$ and $0\leq j \leq n$. We can see how each rotation is separated off as such 

\begin{center}
\begin{tabular}{ c | c c c c c }
$\textbf{A}(t,j)$ & $t=1$ & $t=2$ & $t=3$ & $t=4$ & $t=5$ \\ 
\hline
           & $\textcolor{blue}{\tt{1}}\tt{1233}$, $\textcolor{blue}{\tt{1}}\tt{1323}$  &   $\tt{3}\textcolor{blue}{\tt{1}}\tt{123}$, $\tt{3}\textcolor{blue}{\tt{1}}\tt{132}$ &  $\tt{33}\textcolor{blue}{\tt{1}}\tt{12}$, $\tt{23}\textcolor{blue}{\tt{1}}\tt{13} $ & $\tt{223}\textcolor{blue}{\tt{1}}\tt{1}$, $\tt{323}\textcolor{blue}{\tt{1}}\tt{1}$ & $\tt{1223}\textcolor{blue}{\tt{1}}$, $\tt{1323}\textcolor{blue}{\tt{1}}$
           \\
$j=1$      & $\textcolor{blue}{\tt{1}}\tt{1332}$, $\textcolor{blue}{\tt{1}}\tt{1333}$                & $\tt{2}\textcolor{blue}{\tt{1}}\tt{133}$, $\tt{3}\textcolor{blue}{\tt{1}}\tt{113}$ & $\tt{33}\textcolor{blue}{\tt{1}}\tt{12}$, $\tt{33}\textcolor{blue}{\tt{1}}\tt{12}$ & $\tt{332}\textcolor{blue}{\tt{1}}\tt{1}$, $\tt{333}\textcolor{blue}{\tt{1}}\tt{1}$ & $\tt{1332}\textcolor{blue}{\tt{1}}$, $\tt{1333}\textcolor{blue}{\tt{1}}$
           \\
           \\
           & $\textcolor{blue}{\tt{12}}\tt{233}$, $\textcolor{blue}{\tt{12}}\tt{223}$ & $\tt{3}\textcolor{blue}{\tt{12}}\tt{13}$, $\tt{3}\textcolor{blue}{\tt{12}}\tt{22}$ & $\tt{33}\textcolor{blue}{\tt{12}}\tt{1}$, $\tt{33}\textcolor{blue}{\tt{12}}\tt{1}$ & $\tt{133}\textcolor{blue}{\tt{12}}$, $\tt{223}\textcolor{blue}{\tt{12}}$ & $\textcolor{blue}{\tt{2}}\tt{133}\textcolor{blue}{\tt{1}}$, $\textcolor{blue}{\tt{2}}\tt{223}\textcolor{blue}{\tt{1}}$  \\ 
 $j=2$ & $\textcolor{blue}{\tt{12}}\tt{232}$, $\textcolor{blue}{\tt{12}}\tt{232}$ & $\tt{2}\textcolor{blue}{\tt{12}}\tt{23}$, $\tt{3}\textcolor{blue}{\tt{12}}\tt{23}$ & $\tt{32}\textcolor{blue}{\tt{12}}\tt{2}$, $\tt{33}\textcolor{blue}{\tt{12}}\tt{2}$ & $\tt{232}\textcolor{blue}{\tt{12}}$, $\tt{233}\textcolor{blue}{\tt{12}}$  & $\textcolor{blue}{\tt{2}}\tt{232}\textcolor{blue}{\tt{1}}$, $\textcolor{blue}{\tt{2}}\tt{233}\textcolor{blue}{\tt{1}}$  
 \\
\end{tabular}
\end{center}

\noindent
For $j \in \{0,3,4,5\}$, $\textbf{A}(t,j) = \emptyset$ for all values of $t$. 

\end{exam}
%============================

\noindent
Clearly $A(t,n) = 0$ for all $1 \leq t \leq n$.  Otherwise for $j<n$, each set $\mathbf{A}(t,j)$ falls into one of the following two cases.
\medskip

%========
\noindent
{\bf Case 1} ($t+j \leq n$).  Each $\bs \in \mathbf{A}(t,j)$ is of the form
$\br\,\blue{\tt{a}_1\tt{a}_2\cdots \tt{a}_j}\,\tt{x}\,\bq$ where:
\squishlisttwo
\item $\br \in \bS_k(t-1,w')$ such that every non-empty suffix is larger than $\alpha$, 
\item $\tt{x} < \tt{a}_{j+1}$,
\item $\bq \in \bS_k(n-t-j,w-w'-\wt(\tt{a}_1\cdots \tt{a}_j\tt{x}))$, and 
\item $\wt(\bs) \geq w$.
\squishend

\noindent   
Recall that the number of possibilities for $\br$ with a fixed weight $w'$ is given by $P(t{-}1,0,w')$. 
For each $\br$ with weight $w'$ and given 
$\tt{x}$, there are $S_k(n-t-j,w-w'-\wt(\tt{a}_1\cdots \tt{a}_j\tt{x}))$ possibilities for $\bq$. Thus, in this case
\[ A(t,j)  = \sum_{\tt{x} = 1}^{\tt{a}_{j+1}-1}
\sum_{w'=t-1}^{k(t-1)} P(t-1,0,w')S_k(n-t-j, w-w'-\wt(\tt{a}_1\cdots \tt{a}_j\tt{x}))
. \]



%========
\noindent
{\bf Case 2} ($t+j > n$).  Each $\bs \in \mathbf{A}(t,j)$ is of the form
$\blue{\tt{a}_{n-t+2}\cdots \tt{a}_{j-1}\tt{a}_{j}}\,\tt{x}\,\br\,\blue{\tt{a}_1\tt{a}_2\cdots \tt{a}_{n-t+1}}$ where:
\squishlisttwo
\item $\tt{x} < \tt{a}_{j+1}$,
\item $\br \in \bS(n-j-1, w-\wt(\tt{a}_1\cdots \tt{a}_j\tt{x}))$, and 
\item every non-empty suffix of $\tt{a}_{n-t+2}\cdots \tt{a}_{j-1}\tt{a}_{j}\,\tt{x}\,\br$ is larger than $\alpha$.\footnote{This follows from the definition of $t$ noting that $|a_{n-t+2}\cdots a_{j-1}a_{j}\,x\,\br| < n$.}
\squishend
%
Let $\bq = \tt{a}_{n-t+2}\cdots \tt{a}_{j-1}\tt{a}_{j}$.  Since  $\bq$ is a substring of the necklace $\alpha$,  any suffix of $\bq$ must be lexicographically greater than or equal to the prefix of $\alpha$ with the same length (by the definition of a necklace).  
Therefore we determine the \emph{longest} suffix of $\bq$ that is equal to the prefix of $\alpha$ with the same length.  Suppose this suffix has length $z$.  This
 implies $\tt{a}_{j-z+1}\cdots \tt{a}_{j-1}\tt{a}_{j} = \tt{a}_1\tt{a}_2\cdots \tt{a}_z$.  This means any suffix of $\bs$ starting from an index less than or equal to $|
\bq|-z$ is larger than $\alpha$.  Consider three sub-cases depending on $\tt{x}$.
%
\squishlist
\item If $\tt{x} < \tt{a}_{z+1}$ then $\tt{a}_{j-z+1}\cdots \tt{a}_{j-1}\tt{a}_{j}\tt{x}$ is smaller than $\alpha$, a contradiction.  
\item If $\tt{x} = \tt{a}_{z+1}$ then $\tt{a}_{j-z+1}\cdots \tt{a}_{z+1}\br \in \mathbf{B}(n-j+z, z+1, w-\wt(\tt{a}_1\cdots \tt{a}_{j-z}))$.
 \item If $\tt{x} > \tt{a}_{z+1}$ then any suffix of $\bs$ starting from an index less than or equal to $|\bq|+1$ will be larger than $\alpha$ by Property~\ref{prop:suffix}.  Thus, 
$\br \in \mathbf{B}(n-j-1,0,w-\wt(\tt{a}_1\cdots \tt{a}_j\tt{x}))$.   
 %$A(t, j)= B(n-j-1,0)$.  
 \squishend
Thus, if $\tt{a}_{j+1} > \tt{a}_{z+1}$
%
\[ A(t,j)  = B(n-j+z, z+1, w-\wt(\tt{a}_1\cdots \tt{a}_{j-z})) + 
\sum_{\tt{x} = \tt{a}_{z+1}+1}^\tt{a_{j+1}-1} B(n-j-1,0,w-\wt(\tt{a}_1\cdots \tt{a}_j\tt{x}) ). \]
%
Otherwise $A(t,j) = 0$.

%======================
\subsection{Analysis}

By applying dynamic programming, all necessary $B(t,j,w)$ and $P(t,j,w)$ can be computed in $\mathcal{O}(n^3k^2)$ time using $\mathcal{O}(n^3k)$.  With these values precomputed,
%\begin{itemize}
    %\item 
    $A(t,j)$ can be computed in $O(nk^2)$ time for {\bf Case 1} and $O(k)$ time for {\bf Case 2}, and thus
 %   \item 
    $T_k(n,w,\alpha)$ can be computed in $O(n^3k^2)$ time.
%\end{itemize}
%
\noindent
This proves Lemma~\ref{lem:T}.


%\subsection{Ranking necklaces and Lyndon words with bounded weight}

%\red{Can leave this out for a preliminary version.}


%====================================


\begin{comment}
To compute $A(t,j)$ the set $\mathbf{A}(t,j)$ can be further partitioned into smaller subsets based on two cases: either (i) $t+j > n$ or (ii) $t+j \leq n$.

\medskip 

\noindent
\textbf{Case (i)}: $t+j > n$. For this case, each string $\bs \in \mathbf{A}(t,j)$ is of the form:
%
\[
\tt{a}_{n-t+2}\cdots \tt{a}_{j}\tt{x}\tt{\sigma} \tt{a}_1\tt{a}_2\cdots \tt{a}_{n-t+1}
\]
%
for some $\tt{x} < \tt{a}_{j+1}$ and $\sigma \in \bS_k(n-j-1,w-\wt(\tt{a}_1\cdots \tt{a}_{j}\tt{x}))$.  Furthermore, every suffix of $\tt{a}_{n-t+2}\cdots \tt{a}_{j}\tt{x}\tt{\sigma}$ must be greater than $\alpha$ by the definition of $t$. 
%
%
To enumerate $A(t,j)$, consider the set $\mathbf{B}(t,j,w)$ which denotes the subset of $\bS_k(t,j,w)$ where the strings have prefix $\tt{a}_1\cdots \tt{a}_j$ such that every suffix of each string is greater than $\alpha$.  Let $B(t,j,w) = |\mathbf{B}(t,j,w)|$. It satisfies the following recurrence for $0 \leq j \leq t$:
%
\[
B_k(t,j,w) = 
\begin{cases}	
	1 & \text{if } t=j=0 \text{ and } w\leq 0;  \\
	0 & \text{if } t = j \text{ and } w>0;  \\
	B_k(t,j+1,w) + \displaystyle{\sum_{\tt{x}=\tt{a}_{j+1}+1}^{\tt{k}}}B_k(t{-}j{-}1,0,w{-} \wt(\tt{a}_1\cdots \tt{a}_j\tt{x})) & \text{if } 0 \leq j < t. \\
\end{cases}
\]
\normalsize



\begin{proof} if the $j+1$'st symbol is equal to $\tt{a}_{j+1}$ then these strings are equivalent to $ \textbf{B}(t,j+1,w)$. However, if the symbol is greater, the $j+1$'st symbol could be any number between $\tt{a}_{j+1}+1$ and $k$ inclusive. So consider some value i such that $ \tt{a}_{j+1} \leq i \leq k$ , such that the $j+1$'st symbol be $i$, we must ensure the last $t-j-1$ digits of the string are in $\textbf{B}(t-j-1,0,w-\wt(\alpha[1,j])-i))$. We iterate and sum up over $i$ to get our desired result. We notice that $B(0,0,w) = 1$ for $w \leq 0$ to account for the empty string, and when $t=j\neq0$ results in $B(t,j,w) =0$ or if $t=j=0$ and $w>0$, then the count is 0.
\end{proof}

Using this set to help us enumerate the set of possible values for $\sigma$, take the largest suffix of $\tt{s}_{n-t+2}\cdots \tt{s}_{j-1}\tt{s}_j$ such that it is equal to the prefix of $\alpha$ with the same length. We need to find the possible strings for $\tt{s}_1\tt{s}_2\cdots \tt{s}_s \tt{x}\tt{\sigma}$ such that every non-empty suffix is greater than $\alpha$ with the added restriction that $\tt{s}_{s+1} < \tt{x} < \tt{a}_{j+1}$.The logic of the case when $0\leq j < t$ in Lemma~\ref{lem:B} can be applied to this one, with the adjustment that the sum goes from $\tt{x} = \tt{s}_{s+1}+1$ to $\tt{a}_{j+1}-1$. Note that when $n=j$, the number of possible strings is 0.

So the cases can be summed up as follows. If $n=j$ or if $\tt{s}_{s+1} \geq \tt{a}_{j+1}$ , then $B_k(t,j,w, \alpha) = 0$. In the case that ($n \neq j$ and $\tt{s}_{s+1} = \tt{a}_{j+1}{-}1$), then the count equates to 
\[
B_k(n-j+s,s+1,w- \wt(\alpha[s+1,j]), \alpha) 
\]

\noindent Otherwise, in the case ($n \neq j$ and $\tt{s}_{s+1} < \tt{a}_{j+1}$), then the count equates to 
\[
 B_k(n-j+s,s+1,w- \wt(\alpha[s+1,j]), \alpha)  + \sum_{i=\tt{s}_{s+1}+1}^{\tt{a}_{j+1}-1}{B}_k(t-j-1,0,w-i - \wt(\alpha[s+1,j]), \alpha)
\]



\noindent \textbf{Case (ii)}: $t+j \leq n$. For this case, each string $\bs \in A(t,j)$ is of the form:
\[
\tt{\sigma} \tt{s}_1\tt{s}_2\cdots \tt{s}_j\tt{x}\tt{\tau}
\]

for some $\sigma \in \Sigma^{t-1}$ with the additional property that every suffix is greater than or equal to $\alpha$, $x <  \tt{a}_{j+1}$ and $\tau \in \Sigma^{n-t-j}$ with no additional constraints. If we are able to calculate the total number of possibilities for $\sigma$ and $x\tau$, it would allow us to enumerate the set in this case. 


This motivates the creation of two sets, the already mentioned $\bS_k(n,w^\uparrow)$. The other is the following. 

For some predetermined necklace $\alpha = \tt{a}_1\tt{a}_2\cdots \tt{a}_n$ in some alphabet $\Sigma_{k}$, let $\mathbf{P}(t,j,w)$ represent the set of strings of length $t$, prefix $\alpha[1,j]$ and 
weight equal to $w$. Let $S_k(n,w) = |$\bS_k(n,w^\uparrow)$| and let $P(t,j,w) = |\mathbf{P}(t,j,w)|$.

%We present the recursive formula for calculating the number of $q-ary$ strings of length $n$ such that the sum of its digits is greater than or equal to a given $w$. Let the set of such strings be denoted $\textbf{S}_k(n,w)$ and the cardinality denoted $S_k(n,w)$ 

% Let $W_k^{\uparrow}(n,w)$ denote the cardinality of $\bS_k(n,w^\uparrow)$


\begin{lemma} \label{lem:count}
    %The number of all $q$-ary strings of length $n$ and maximum weight $w$, can be represented by the following recursive function, $S(n,w)$.
    For $k\geq 2$, $n\geq 0$, and an integer weight $w$: 
    \[
    S_k(n,w) = 
    \begin{cases}
        0 & \mbox{if $(n=0$ \text{ and } $w>0)$ \text{ or } $w > kn$;}\\
        k^n & \mbox{if $w\leq 0$;}\\
        \displaystyle{\sum_{i = 1}^{k}S_k(n-1,w-i)} & \text{otherwise.} \\
    \end{cases}
    \]
\end{lemma}
%
% \begin{proof}
% %
% We notice that if $n=0$ and $w>0$, its impossible to create an empty string that has a weight bigger than $0$, and if $w=0$ we know that it is possible to create a string of length $0$ and weight $0$. We note that if $w=0$, then the possible number of strings is $k^n$. As for the recurrence, in any string with $n>1$, we notice that we can fix the string at the front with any of the numbers in $\Sigma_k$. The respective remaining weights would be $\{w,w-1,\ldots,w-k\}$, and so can sum all of those possibilities, we get the recursive function we mention above. We also notice that it is possible for $w-k<0$, in this case it is the exact same as if the the weight was 0 instead of $w-k$, so the number of possible strings is still $k^n$.
% %
% \end{proof}

Recall that $\alpha = \tt{a}_1\tt{a}_2\cdots \tt{a}_n$ is a fixed necklace.  \red{what is $t$}.


\begin{lemma} For $k \geq 2$, $n \gea 0$, and $0 \leq j \leq t $.
\[
\small 
{P}_k(t,j,w) = 
\begin{cases}	
	1 & \text{if } t=j=w=0 \\
	0 & \text{if } t = j\text{ and } w\neq0 \\
	{P}_k(t,j+1,w) + \displaystyle{\sum_{i=a_{j+1}+1}^{k}P_k(t-j-1,0,w - \wt(\alpha[1,j]) - i)} & \text{if } 0 \leq j < t \\
\end{cases}
\]
\normalsize
\end{lemma}


Using this to enumerate the possibilities for $\sigma$ and $x\tau$, the weight gets divided between $\tt{\sigma}$ and $\tt{x}\tt{\tau}$, so for all possible values of the weight in $\tt{\sigma}$, we need to ensure that the weight of $\tt{x}\tt{\tau}$ and $\tt{\sigma}$ sum up to greater than $w$. Thus we iterate from $q=1$ to $q=k$ and have the weight for $\tt{\sigma}$ be equal to $q$ and the weight for $x\tau$ be greater than or equal to $w-q$. We will call $w - \wt(\alpha[1,j]) = w'$ for simplicity. 

As for the $\tt{x}\tt{\tau}$, this is essentially equivalent to Lemma~\ref{lem:count}, but we note that the first digit is not between $1$ and $k$, rather it is between $1$ and $\tt{s}_{j+1}-1$ inclusive. Note that if $\tt{s}_{j+1} = 1$, then there are 0 possible strings, in any other case the following holds.
\[
	\sum_{i=1}^{\tt{s}_{j+1}-1}S(n-1,w-i)
\]

\noindent So now we sum up over all possible values of q, we get the following 
\[
	\sum_{q=0}^{k(t-1)} P_k(t-1,0,q, \alpha)(\sum_{i=1}^{\tt{s}_{j+1}-1}S(n-1,w-q-i))
\]
\normalsize

\noindent
By applying dynamic programming, 

\noindent - $W_k^{\uparrow}(n,w)$ can be computed in $O(n^2k^2)$ time using $O(n^2k)$ space.

\noindent - $B(t,j,w)$ can be computed in $O(n^3k^2)$ time using $O(n^3k)$ space.

\noindent - $P(t,j,w)$ can be computed in $O(n^3k^2)$ time using $O(n^3k)$ space.

\noindent - $A(t,j,w)$ can be computed in $O(k)$ time using $O(1)$ space.

\noindent Therefore, $T(n,w)$ can be computed in $O(n^3k^2)$ time, proving Lemma~\ref{lem:T}.



% \red{ Provide a unified statement that counts $A$ at the end}


% \red{Wrap up the computation of $T$ and the running time / space and say it proves the Lemma~\ref{lem:T}.}

% In either case, if the tables of $B$,$P$ and $S$ are complete, the maximum computational time is of order $O(n^2k)$, 

% Thus, in conclusion, T can be computed in $O(n^3k^2)$ time with $O(n^3k)$ space, proving Lemma~\ref{lem:T}.

% \begin{algorithm}
% \small
% \caption{Computing $T_k(n, w, \alpha)$ for a given $k$-ary string $\alpha = \tt{a_1 a_2 \cdots a_n}$}.
% \begin{algorithmic}[1]


% \Function{T}{$\alpha,n,w$}
%     \Statex $q$ is given a predetermined value.
%     \State $\alpha \gets \textsc{LargestNecklace}(\alpha, n)$
%     \State
%     \State $total \gets 0$

%     \State \textcolor{blue}{\textbf{Precompute} $S(n,w)$ using dynamic programming}
%     \For{i \textbf{from} 1 to kn}
%        \State  $S(0,i) \gets 0$
%     \EndFor
%     \For{i \textbf{from} 0 to n}
%        \State  $S(i,0) \gets i^k$
%     \EndFor
%     \For{i \textbf{from} 1 to n}
%         \For{j \textbf{from} 1 to kn}
%             \For{$q \textbf{ from } 0 \textbf{ to } k$}
%                 \State $total \gets total + S(n-1,w-q)$
%             \EndFor
%             \State $S(n,w) \gets total$
%         \EndFor
%     \EndFor

%     \State

%     \State $total \gets 0$

    
%     \State \textcolor{blue}{\textbf{Precompute} $P_\alpha(t, j, w)$ using dynamic programming}
%     \State $P_\alpha(0,0,0) \gets 1$
%     \For{$t \textbf{ from } 1$ to $n$}
%         \For{$w \textbf{ from } 0$ to $n$}
%             \State $P_\alpha(t, t, w) \gets 0$
%             \For{$j \gets t-1$ down to $0$}
%                 \State $total \gets P(t,j+1,w)$ 
%                 \For{$ q\textbf{ from } a_{j+1} +1 \textbf{ to } k $}
%                     \If{$w-k-W(\alpha[1,j]) < 0$ }
%                         \State $\textbf{break}$
%                     \EndIf\If{$w-k-W(\alpha[1,j]) < 0$ }
%                         \State $\textbf{break}$
%                     \EndIf
%                     \State $total \gets total + P(t-j-1,0,w-k-W(\alpha[1,j]))$
%                 \EndFor
%                 \State $P(t,j,w) \gets total$
%             \EndFor
%         \EndFor
%     \EndFor
%     \State 

%     \State $total \gets 0$
    
%     \State \textcolor{blue}{\textbf{Precompute} $B_\alpha(t, j, w)$ using dynamic programming}
%     \State $B_\alpha(0,0,0) \gets 1$
%     \For{$t \textbf{ from } 1$ to $n$}
%         \For{$w \textbf{ from } 0$ to $n$}
%             \State $B_\alpha(t, t, w) \gets 0$
%             \For{$j \gets t-1$ down to $0$}
%                 \State $total \gets B(t,j+1,w)$ 
%                 \For{$ q\textbf{ from } a_{j+1} +1 \textbf{ to } k $}
%                     \If{$w-k-W(\alpha[1,j]) < 0$ }
%                         \State $total \gets total + B(t-j-1,0,0)$
%                     \Else
%                         \State $total \gets total + B(t-j-1,0,w-k-W(\alpha[1,j]))$
%                     \EndIf
%                 \EndFor
%                 \State $B(t,j,w) \gets total$
%             \EndFor
%         \EndFor
%     \EndFor
        
% \EndFunction
        
 
% \end{algorithmic}
% \end{algorithm}


% \normalsize

\begin{comment}
%=======================================
%=======================================

\end{comment}


%Future add on ranking problems?
%\begin{enumerate}
 %   \item Rank necklaces with fixed weight
  %  \item Rank necklaces with fixed content
   % \item Rank necklaces with bounded weight (lower) (for subsets)
    %\item What is the lex largest necklace for given content/weight? Matthew work follow up.
%\end{enumerate}

\begin{comment}
%=======================================
\section{Summary and future work} \label{sec:future}

In this paper we detailed the first efficient algorithms to decode bounded weight de Bruijn sequences, universal cycles for $t$-subsets, and universal cycles for $t$-multisets.  Implementations of these algorithms can be found at \url{debruijnsequence.org}.  An interesting line of future research is to develop efficient decoding algorithms for the following universal cycles:
\begin{enumerate}
    \item any fixed-weight de Bruijn sequence~\cite{fixedweight2}, which corresponds to a universal cycle for $t$-subsets using the shorthand binary representation,
    \item the Grandmama de Bruijn sequence~\cite{grandma2},
    \item the de Bruijn sequence derived from lexicographic compositions~\cite{lexcomp},
    \item the prefer-same de Bruijn sequence~\cite{pref-same}
    \item the prefer-opposite de Bruijn sequence~\cite{pref-opp},
    \item any cut-down de Bruijn sequence~\cite{etzion-cut,cutdown,nellore},
    \item any universal cycle for multiset permutations~\cite{coolest}.

\end{enumerate}

We also note that the results from Section~\ref{sec:rank} can be applied to efficiently rank necklaces and Lyndon words with a lower bound on their weight.  

\end{comment}

%=======================================
%=========================
\section{Acknowledgment}
Joe Sawada (grant RGPIN-2025-03961) gratefully acknowledges research support from the Natural Sciences and Engineering Research Council of Canada (NSERC).



%------------------------
%\begin{exam} 
%Consider $n=3$, $k=4$, and $w=9$.  Since 
%\[  \bN^{\uparrow}_4(3,9)  = \tt{144,~ 234,~ 243,~ %244,~ 333,~ 334,~ 344,~ 444}, \]
%we have
%\[ \mathcal{W}_4^{\uparrow}(3,9) = 
%\tt{144\cdot
%234\cdot
%243\cdot
%244\cdot
%3\cdot
%334\cdot
%344\cdot
%4}.\]
%
%Consider $\alpha_1 = 312$. In $\mathcal{G}_4(3)$, $\alpha_1$ is found as a substring of $113 \cdot 122$; however, in $\mathcal{DB}_4(3,6)$ it is found as a substring of $033 \cdot 123$. Similarly,  $\alpha_2 = 303$ is found as a substring of $023 \cdot 031$ in the $\mathcal{G}_4(3)$, but is found in the wraparound of $\mathcal{DB}_4(3,6)$.    
%\end{exam}
%------------------------



%-----------  BIBLIOGRAPHY  ---------------


\bibliographystyle{alphaabbr}
\bibliography{refs}
\normalsize
\begin{comment}
\appendix

\section{CODE}

\small

\begin{code}
#include <stdio.h>
#include <stdlib.h>
#include <string.h>
#include <math.h>

//Table for S
#define MAXN 100
#define MAXW 100

int tableS[MAXN][MAXW];

//Table for B
#define MAXBT 100
#define MAXBJ 100
#define MAXBW 100

int tableB[MAXBT][MAXBJ][MAXBW];

//Table for P
#define MAXPT 100
#define MAXPJ 100
#define MAXPW 100

int tableP[MAXPT][MAXPJ][MAXPW];



// --------------------------------------------------------------------------
// ----------------------------HELPER FUNCTIONS------------------------------
// --------------------------------------------------------------------------


// USED FOR A

int sumWeightFirstJ(int alpha[], int j) {
    int weight = 0;
    for (int i = 0; i < j; i++) {
        weight = weight + alpha[i];
    }
    return weight;
}


// Computes n choose k (binomial coefficient)
unsigned long long choose(int n, int k) {
    if (k < 0 || k > n) return 0;
    if (k == 0 || k == n) return 1;
    if (k > n - k) k = n - k; // Use symmetry
    unsigned long long result = 1;
    for (int i = 1; i <= k; ++i) {
        result = result * (n - i + 1) / i;
    }
    return result;
}

long long int numStringsWithFirstDigitLimit(int w, int n, int q, int p) {
    long long int total = choose(w+n, n);
    for (int i=0;w-p-i*q>=0;i++){
        total += choose(n-1,i)*choose(w-p-i*q+n,n)*pow(-1,i+1);
    }
    for (int i=1;w-i*q>=0;i++){
        total += choose(n-1,i)*choose(w-i*q+n,n)*pow(-1,i);
    }
    return total;
}

int checkPreSuf(int alpha[],int pre1,int pre2,int suf1){
    for (int i=0;i<pre2-pre1;i++){
        if (alpha[pre1+i] != alpha[suf1+i]){
            //if there are two indexes that are not equal we return a false value. 
            return 0;
        }
    }
    //this only goes through if all the indexes are equal
    //note if pre2 == pre1 then this will also return true. 
    return 1;
}


int findLargestS(int alpha[], int j, int t, int n){
    int largestS = 0;
    for (int i=1;i<=j-(n-t+1);i++){
        //looks through the first j-(n-t+1) elements of alpha.
        if (checkPreSuf(alpha,j-i,j,0)){
            //looks for longest suffix
            largestS = i;
        }
    }
    return largestS;
}


//USED FOR PHI

int gcd(int a, int b) {
    while (b != 0) {
        int temp = b;
        b = a % b;
        a = temp;
    }
    return a;
}



//USED FOR N 

int isGreater(int *arr, int *alpha, int n) {
    for (int i = 0; i < n; ++i) {
        if (arr[i] > alpha[i]) return 1;
        if (arr[i] < alpha[i]) return 0;
    }
    return 0; // equal
}



// Helper function to check if array is lexicographically minimal among its rotations
int isNecklace(int *arr, int n) {
    for (int shift = 1; shift < n; ++shift) {
        int is_smaller = 0;
        for (int i = 0; i < n; ++i) {
            int a = arr[i];
            int b = arr[(i + shift) % n];
            if (a < b) break;
            if (a > b) { is_smaller = 1; break; }
        }
        if (is_smaller) return 0;
    }
    return 1;
}

// Finds the next largest necklace above alpha (lex order)
int nextLargestNecklace(int *alpha, int n, int q) {
    int *arr;
    arr = malloc(n * sizeof(int));
    memcpy(arr, alpha, n * sizeof(int));
    while (1) {
        // Increment lexicographically
        int i = n - 1;
        while (i >= 0 && arr[i] == q - 1) {
            arr[i] = 0;
            i--;
        }
        if (i < 0) return 0; // No larger string exists
        arr[i]++;
        // Check if it's a necklace and greater than alpha
        if (isNecklace(arr, n) && isGreater(arr, alpha, n)) {
            memcpy(alpha, arr, n * sizeof(int));
            free(arr);
            return 1;
        }
    }
    
}





// --------------------------------------------------------------------------
// ----------------------------END HELPER FUNCTIONS--------------------------
// --------------------------------------------------------------------------





//FUNCTION TO CALCULATE B(t,j,w,q)
int B(int t, int j, int w, int alpha[], int q) {
    int total=0;
    int WinJ = 0;
    if ( w<0 || t<0 || j<0) {
        return 0;
    }
    if (tableB[t][j][w] == -1){
        for (int i = 0; i < j; i++) {
            WinJ = WinJ + alpha[i];
        }
        total += B(t, j + 1, w, alpha, q);
        for (int k = alpha[j] +1; k < q; k++) {
            total += B(t - j - 1, 0, w - k - WinJ, alpha, q);
        }
        tableB[t][j][w] = total;
    } else {
        total = tableB[t][j][w];
    }
    return tableB[t][j][w];
}

//FUNCTION TO CALCULATE P(t,j,w,q)
int P(int t, int j, int w, int alpha[], int q) {
    int total=0;
    int WinJ = 0; //Weight within the range of 1 to J
    if (t<0 || j<0 || w<0 || w> t*(q-1)) {
        return 0;
    }
    if (tableP[t][j][w] == -1){
        for (int i = 0; i < j; i++) {
            WinJ = WinJ + alpha[i] ;
        }
        total += P(t, j + 1, w, alpha, q);
        for (int k = alpha[j] +1;k < q; k++) {
            total += P(t - j - 1, 0, w - k - WinJ, alpha, q);
        }
        tableP[t][j][w] = total; // Store the result in the table
    } else {
        total = tableP[t][j][w];
    }
    return tableP[t][j][w];
}


//FUNCTION TO CALCULATE S(n,w,q)
int S(int n, int w, int q){
    int buffer = 0;
    if (tableS[n][w] != -1) {
        return tableS[n][w];
    }else{
        for (int i=0;i<q && w-i>=0;i++){
            buffer+= S(n-1, w-i, q);
            tableS[n][w] = buffer;
        }
        return tableS[n][w];
    }
}


//FUNCTION TO CALCULATE the alternate S(n,w,q) in case 1 of A. 
int xS(int n, int w, int q, int j){
    int total = 0;
    for (int i=0;i<j && i <= w;i++){
        total += S(n-1, w-i, q);
        tableS[n][w] = total; // Store the result in the array
    }
    return total;
}


//Function to calculate A(t,j,w,q) 
int A(int n, int t, int j, int w, int q, int alpha[]){
    int total = 0;
    int s;
    int w1 = w - sumWeightFirstJ(alpha, j);
    // Initialize the array with -1
    for (int i = 0; i < MAXN; i++) {
        for (int j = 0; j < MAXW; j++) {
            tableS[i][j] = -1;
            if (i == 0) {
                tableS[i][j] = 1; // There's one way to form a string of length 0
            }
            if (j == 0) {
                tableS[i][j] = 1; // There's one way to form a string of weight 0
            }
        }
    }
    if (n==j && sumWeightFirstJ(alpha, j) <= w){
        return 1;
    }
    if (t+j  <= n){
        for (int k=0;k<=w1;k++){ 
            total += P(t-1,0,k,alpha,q)*(xS(n-t-j+1,w1-k,q,alpha[j]));
        }
    }
    else {
        s = findLargestS(alpha,j,t,n);
        if (alpha[s] < alpha[j]){
            total += B(n-j+s,s+1,w1+sumWeightFirstJ(alpha,s),alpha,q);
        }
        for (int i=alpha[s]+1;i<alpha[j];i++){
            total += B(n-j-1,0, w1 + sumWeightFirstJ(alpha,s)-i,alpha, q);
        }
    }
    return total;
}

//Function to caulcate T(n,w,q) 
int T(int n, int w, int q, int alpha[]){
    //Setting up Table for P
    memset(tableP, -1, sizeof(tableP));
    for (int t=1;t<MAXPT;t++){
        for (int i=0;i<MAXPJ;i++){
            tableP[t][t][i] = 0;
        }
    }
    tableP[0][0][0] = 1; // Base case: P(0, 0, 0) = 1

    //Setting up Table for B
    memset(tableB, -1, sizeof(tableB));
    for (int t=1;t<MAXPT;t++){
        for (int i=0;i<MAXPJ;i++){
            tableB[t][t][i] = 0;
        }
    }
    for (int w=0;w<MAXBW;w++){
        tableB[0][0][w] = 1; // Base case: B(0, 0, w) = 1
    }


    int totalStrings = 0;
    for (int t = 1; t <= n; t++) {
        for (int j = 0; j <= n-1; j++) {
            totalStrings += A(n, t, j, w, q, alpha);
        }
    }
    return totalStrings;
}

//Function to calculate phi(n)
int phi(int n){
    int value = 0;
    if (n == 1) {
        return 1; // phi(1) = 1
    }
    for (int i = 1; i < n; i++) {
        if (gcd(n, i) == 1) {
            value++;
        }
    }
    return value;
}


//Function to calculate N(n,w,q) 
int N(int n, int w, int q, int alpha[]){
    int *result;
    int total = 0;
    for (int i=1;i<=n;i++){
        if (n%i == 0){
            result = malloc(n/i * sizeof(int));
            for (int j=0;j<n/i;j++){
                result[j] = alpha[j];
            }
            nextLargestNecklace(result, n/i, q);
            total += T(n/i, w/i, q, result) * phi(i);
            free(result);
        }
    }
    return total/n;
}

    
\end{code}

\end{comment}






\begin{comment}

Observe that the Granddaddy
% 
\[ 
\mathcal{G}_4(3) = 
\red{0~
001~
002~
003~
011~
012~
013~
021~
022~
023~
031~
032~}
\underline{033}~
\red{1~
112~
113~
122~}
\underline{123~
132~
133~
2~
223~
233~
3}\]
\normalsize

0 0 0
0 0 1
0 0 2
0 0 3
0 1 1
0 1 2
0 1 3
0 2 1
0 2 2
0 2 3
0 3 1
0 3 2
0 3 3
    1 1 1
    1 1 2
    1 1 3
    1 2 2
1 2 3
1 3 2
1 3 3
2 2 2
2 2 3
2 3 3
3 3 3


Consider all necklaces of length $n=4$ over alphabet of size $k=4$.  Suppose we bound the weight at 4.  Which necklaces do we need to remove?

    0 0 0 0
    0 0 0 1
    0 0 0 2
    0 0 0 3
    0 0 1 1
    0 0 1 2
0 0 1 3
    0 0 2 1
0 0 2 2
0 0 2 3
0 0 3 1
0 0 3 2
0 0 3 3
    0 1 0 1
    0 1 0 2
0 1 0 3
    0 1 1 1
0 1 1 2
0 1 1 3
0 1 2 1
0 1 2 2
0 1 2 3
0 1 3 1
0 1 3 2
0 1 3 3
0 2 0 2
0 2 0 3
0 2 1 1
0 2 1 2
0 2 1 3
0 2 2 1
0 2 2 2
0 2 2 3
0 2 3 1
0 2 3 2
0 2 3 3
0 3 0 3
0 3 1 1
0 3 1 2
0 3 1 3
0 3 2 1
0 3 2 2
0 3 2 3
0 3 3 1
0 3 3 2
0 3 3 3
1 1 1 1
1 1 1 2
1 1 1 3
1 1 2 2
1 1 2 3
1 1 3 2
1 1 3 3
1 2 1 2
1 2 1 3
1 2 2 2
1 2 2 3
1 2 3 2
1 2 3 3
1 3 1 3
1 3 2 2
1 3 2 3
1 3 3 2
1 3 3 3
2 2 2 2
2 2 2 3
2 2 3 3
2 3 2 3
2 3 3 3
3 3 3 3

\section{Proofs} \label{sec:proofs}

In this section we prove our two main technical results required to decode the universal cycle $\mathcal{G}_k(n,w^\uparrow)$.

%=====================================================================
\subsection{Proof of Lemma~\ref{lem:alpha2}} \label{sec:proof1}

Assume $n,k > 1$ and $w \geq k$.
Let $\bs=\bp\bq$ be a string in $\bS_k(n,w^\uparrow)$ where $\bq$ is the longest suffix of $\alpha$ such that $\bq\bp$ is a necklace.  Assume $\alpha \neq \tt{k}^n$ and $\alpha$ is not found in the wraparound of $\mathcal{G}_k(n,w^\uparrow)$. 
%
We must demonstrate that:
\begin{enumerate}
    \item[(a)] there exists consecutive necklaces $\beta_1$, $\beta_2$, and $\beta_3$ in $\bN_k(n,w^\uparrow)$ (considered cyclically) such that $\bp$ is a suffix of $\beta_1$ and $\bq$ is a prefix of $\beta_2$,

    \item[(b)] $\alpha$ is a substring of $\ap(\beta_1)\ap(\beta_2)\ap(\beta_3)$, 
    \item[(c)] $\beta_1$, $\beta_2$, and $\beta_3$ are unique,
    and 
    \item[(d)] $\beta_2$ can be determined in $\mathcal{O}(n^2)$ time using $\mathcal{O}(n)$ space.
\end{enumerate}


\noindent
{\bf Item (a).}~~Recall from Step 1 of Section~\ref{sec:granddaddy}, that there exists unique consecutive necklaces $\beta'_1$ and $\beta'_2$ in $\bN_k(n)$ (considered cyclically) such that $\bp$ is a suffix of $\beta'_1$ and $\bq$ is a prefix of $\beta'_2$. Let $\beta_1$ denote the  \emph{largest} necklace in $\bN_k(n,w^\uparrow)$ that is less than or equal to $\beta'_1$. Let $\beta_2$ denote the \emph{smallest} necklace in $\bN_k(n,w^\uparrow)$ that is greater than or equal to $\beta'_2$. 
\begin{quote}
{\bf Claim:} $\beta_1$ and $\beta_2$ exist and furthermore $\beta_1$ has suffix $\bp$ and $\beta_2$ has prefix $\bq$. 
\end{quote}
Clearly $\beta_2$ immediately follows $\beta_1$ in the ordering $\bN_k(n,w^\uparrow)$. Trivially, let $\beta_3$ be the necklace immediately following $\beta_2$.  We now validate the claim.
\begin{itemize}
    \item Consider $\beta_1$.  If $\beta_1 = \beta'_1$, then we are done. Otherwise, $\beta_1$ has weight less than the necklace $\bq\bp$. \red{Add remaining argument}.


    \item Consider $\beta_2$.  If $\beta_2 = \beta'_2$, then we are done.  Otherwise, $\beta_2$ has weight less than the necklace $\bq\bp$. Thus, $\beta_2$ is lexicographically smaller than $\bq \tt{k}^{|\bp|}$, which is a necklace from repeated application of Property~\ref{prop:lastnon}. This implies that  $\beta_2$ exists and moreover, it must have prefix $\bq$. In particular, it can easily be verified that $\beta_2$ is obtained from $\beta'_2$ by repeated application of Property~\ref{prop:lastnon} until it has weight $w$.
    
\end{itemize}

\smallskip
%==============
\noindent
{\bf Item (b).}  We must handle the cases when either $\beta_1$ or $\beta_2$ are periodic.  
From Property~\ref{prop:periodic}, $\bq$ is a prefix of $\ap(\beta_2)\ap(\beta_3)$.  If $\beta_1$ is periodic, suppose that $|\bp| > |\ap(\beta_1)|$.  Then from Property~\ref{prop:periodic}, it must be that $\alpha$ is a rotation of the necklace $\beta_1$. But this string must exist as a substring of $\


\smallskip
%==============
\noindent
{\bf Item (c).}~~Observe that $\beta_1$, $\beta_2$, $\beta_3$ must be unique since $\mathcal{G}_k(n,w^\uparrow)$ is a universal cycle. 

\smallskip
\noindent
{\bf Item (d).}~~Recall from Section~\ref{sec:granddaddy} that $\beta'_2$ can be determined in $\mathcal{O}(n^2)$ time using $\mathcal{O}(n)$ space. Furthermore, as described in Item (a), $\beta_2$ can be easily computed from $\beta'_2$ in $O(n)$ time using $O(n)$ space since each symbol in the updated suffix can be computed in constant time.


\red{\hrule}

\bigskip

% Step zero: define the problem

As stated in Section \ref{sec:granddaddy}, by applying the algorithm outlined in~\cite{kociumaka,ranking} we can determine the location of $\beta_2$, the necklace containing the suffix $\bq$ of any string $\alpha = \bp \bq$ by considering the sequence of necklaces $\bN_k(n)$. To find the position of $\alpha$ in $\bN_k(n,w^\uparrow)$, we utilize this algorithm.

\begin{remark} \label{rem:state}
Let $n,k \geq 1$ and $w \geq 0$. Let $\alpha = \bp\bq$ be a string in $\Sigma_k^n$ with weight at least $w$ such that $\bq$ is the longest suffix of $\gamma$ such that $\bq\bp$ is a necklace. As $\alpha$ is a string in $\Sigma_k^n$, the algorithm outlined in Section \ref{sec:granddaddy} can be used to determine the position of $\beta_2$ in $\bN_k(n)$ in $\mathcal{O}(n^2)$ time using $\mathcal{O}(n)$ space.
\end{remark}

To apply this to $\bN_k(n,w^\uparrow)$ we must also note the following lemmas regarding $\bN_k(n)$, which allows us to prove the position of $\alpha$ in $\bN_k(n,w^\uparrow)$, assuming its position is known in $\bN_k(n)$. 

% Step one: lay out our tools

\begin{lemma} \label{lem:max}
    Let $\beta = \tt{b}_1 \tt{b}_2 \cdots \tt{b}_{i-1} \cdot \tt{x} \cdot \tt{k}^{n-i}$ be a necklace in $\bN_k(n)$ and let $\tt{x}<k$. There exists a necklace $\beta_2 = \tt{b}_1 \tt{b}_2 \cdots \tt{b}_{i-1} \cdot \tt{(x+1)} \cdot \tt{k}^{n-i}$. In other words, increasing the last non-$k$ character of a necklace always produces a valid necklace.

    % Let $\beta$ be a necklace in $\bN_k(n)$ and let $\beta = \bp \bq$, where $\bq$ has length $i$. If the suffix $\bq$ is replaced with a maximal suffix, the resultant string $\bp \cdot \tt{k}^i$ is also a valid necklace in $\bN_k(n)$. If $\beta \neq \bp \cdot \tt{k}^i$, then $\bp \cdot \tt{k}^i$ occurs after $\beta$ lexicographically.
\end{lemma}

\begin{proof}

% Suppose that $\beta = \tt{b}_1 \tt{b}_2 \cdots \tt{b}_n$ is a string of the form $\gamma \cdot \tt{k}^{n-m} \in \Sigma_k^n$. If $\beta$ is not in $\bN_k(n)$, then there exists some (circular) substring $\alpha = \tt{a}_1 \tt{a}_2 \cdots \tt{a}_m$ of $\beta$ where $\alpha_i = \beta_i$ for all $i < m$ and $\alpha_m < \beta_m$. If $\alpha$ begins in the suffix $\tt{k}^{n-m}$, then $\alpha \leq \beta$ when considered cyclically. If $\alpha$ ends in the suffix $\tt{k}^{n-m}$, then $\alpha_m \geq \beta_m$. Thus, $\alpha$ resides entirely within $\gamma$. Therefore, if $\gamma \cdot \tt{k}^{n-m}$ is not a necklace, then every string where $\gamma$ is a prefix is also not a necklace. By the contrapositive, if $\bp \bq$ is a necklace in $\bN_k(n)$, then $\gamma \cdot \tt{k}^{n-m}$ is also a necklace.

Let $\beta = \tt{b}_1 \tt{b}_2 \cdots \tt{b}_{i-1} \cdot \tt{x} \cdot \tt{k}^{n-i}$ be a necklace in $\bN_k(n)$ and let $\tt{x}<k$. By the definition of a necklace, any prefix of $\beta$ must be at most lexicographically equal to any other substring of $\beta$ considered cyclically. This also implies $\tt{b}_1 \leq \tt{x}$ and $\tt{b}_1 < k$. 

Suppose there is another string of the form $\beta_2 = \tt{b}_1 \tt{b}_2 \cdots \tt{b}_{i-1} \cdot \tt{(x+1)} \cdot \tt{k}^{n-i}$. Since $\tt{b}_1< \tt{(x+1)}$ and $\tt{b}_1<k$, there cannot be any substring of $\beta_2$ starting after index $i-1$ that is smaller than its prefix. Likewise, any substring starting between index $2$ and $i-1$ would either be unaffected or increased at a lesser index than the prefix of $\beta_2$. Therefore, as no smaller substring than the prefix is created, $\beta_2$ is also a necklace.

\end{proof}

\begin{exam}
Let $n=4$, $k=3$, with $\beta = \tt{1122}$ in $\bN^{\uparrow}_3(4)$. Consider all possible maximal suffixes $\gamma_i$, where $i$ is the length of $\gamma_i$. Using Lemma \ref{lem:max}, we know that the following necklaces $\beta_i$ also exist after $\beta$ in $\bN^{\uparrow}_3(4)$:

\begin{center}
\begin{tabular}{l c}
    $\bp=\tt{\textcolor{red}{3}}$ & $\beta_1 = \tt{112\textcolor{red}{3}}$ \\
    $\bq=\tt{\textcolor{red}{33}}$ & $\beta_2 = \tt{11\textcolor{red}{33}}$ \\
    $\gamma_3=\tt{\textcolor{red}{333}}$ & $\beta_3 = \tt{1\textcolor{red}{333}}$ \\
    $\gamma_4=\tt{\textcolor{red}{3333}}$ & $\beta_4 = \tt{\textcolor{red}{3333}}$
\end{tabular}
\end{center}
\end{exam}

By Lemma \ref{lem:max}, for any necklace $\beta$ in $\bN_k(n)$ where $\beta = \bp \bq$ and $\bq$ has length $i$, we can replace the suffix $\bq$ with a suffix of the form $\tt{k}^i$ by repeatedly incrementing the final character. If $\beta \neq \bp \cdot \tt{k}^i$, then $\bp \cdot \tt{k}^i$ occurs after $\beta$ lexicographically. We call a suffix of the form $\tt{k}^i$ a \defo{maximal suffix} over $\Sigma_k$.

Additionally, if two sequences begin with the same substring $\tt{a}_1 \tt{a}_2 \cdots \tt{a}_{i-1}$ until some character at index $i$, then those strings share a \defo{common prefix} of length $i-1$. The index where the two strings differ, $i$, is called the \defo{increasing index}. Assuming two strings have the same length, the string with the larger character at the increasing index must be the lexicographically larger of the two. Through these definitions, we can utilize the following lemma from~\cite{grandmama}.

\begin{lemma} \label{lem:lexorder}
    Let $\beta_1$, $\beta_2$ be consecutive necklaces in the sequence of necklaces $\bN_k(n)$. There exists an index $1 \leq i\leq n$ such that
    \[
    \beta_1 = \tt{a}_1 \tt{a}_2\cdots \tt{a}_{i-1} \cdot \tt{x} \cdot \tt{a}_{i+1} \tt{a}_{i+2} \cdots \tt{a}_n
    \]\[
    \beta_2 = \tt{a}_1 \tt{a}_2\cdots \tt{a}_{i-1} \cdot \tt{y} \cdot \tt{b}_{i+1} \tt{b}_{i+2} \cdots \tt{b}_n
    \]
    where $\tt{x},\tt{y}\in \Sigma_k$ with $\tt{x}<\tt{y}$. $\beta_1$ and $\beta_2$ share the common prefix $\tt{a}_1 \tt{a}_2\cdots \tt{a}_{i-1}$, and the increasing index between them is $i$.
    
    % We refer to $\tt{a}_1 \tt{a}_2\cdots \tt{a}_{i-1}$ as the \defo{common prefix} of $\beta_1$ and $\beta_2$, and $i$ as the \defo{increasing index} of $\beta_1$ and $\beta_2$.
\end{lemma}

Given these two lemmas, let us consider the suffix of $\beta_1$, for two necklaces $\beta_1$ and $\beta_2$ that are consecutive in $\bN_k(n)$. Let us represent $\beta_1$ as $\tt{a}_1 \tt{a}_2\cdots \tt{a}_n$. Suppose that the suffix of $\beta_1$ following the increasing index $i$ is not maximal. If so, then we could use Lemma \ref{lem:max} to construct a new necklace of the form $\tt{a}_1 \tt{a}_2\cdots \tt{a}_{i-1} \cdot \tt{x} \cdot \tt{k}^{n-i}$ that is larger than $\beta_1$ but smaller than $\beta_2$ lexicographically. As $\beta_1$ and $\beta_2$ are consecutive, this is a contradiction. This leads us to the following corollary:

\begin{corollary} \label{cor:max2}
    For two necklaces $\beta_1$ and $\beta_2$ to be consecutive in $\bN_k(n)$, the suffix after the increasing index of $\beta_1$ must be a maximal suffix.
\end{corollary}

% Step two: use these tools to prove via two cases

Knowing the position of $\alpha$ in the sequence $\bN_k(n)$, and having established the key attributes of this sequence, we can now consider where $\alpha$ may appear in $\bN_k(n,w^\uparrow)$. There are several possibilities regarding the necklaces $\beta_1$ and $\beta_2$ that contain $\alpha$. 

\begin{itemize}
    \item If both necklaces have a weight of at least $w$, then $\beta_2$ still contains $\bq$ in $\bN_k(n,w^\uparrow)$, and we can use the previously stated algorithm to return $\beta_2$ in $\mathcal{O}(n^2)$ time.
    \item If either $\beta_1$ or $\beta_2$ is not in the sequence $\bN_k(n,w^\uparrow)$, then it is possible that the position of $\alpha$, which can only be found where one necklace ends with the suffix $\bp$ and the following necklace begins with the prefix $\bq$, has shifted to somewhere else within the sequence.
\end{itemize}

\begin{exam}  Let $n=3$, $k=4$, and $w=9$.  Then 

\[  \bN^{\uparrow}_4(3,9)  = \tt{144,~ 234,~ 243,~ 244,~ 333,~ 334,~ 344,~ 444}, \]

\begin{itemize}
    \item In $\bN^{\uparrow}_4(3)$, $\alpha = \tt{423}$ is found as a substring of $\beta_1\beta_2 = \tt{224} \cdot \tt{233}$ 
    \item In the corresponding sequence $\bN^{\uparrow}_4(3,9)$ it is found as a substring of $\tt{144} \cdot \tt{234}$. 
\end{itemize}

So, it is possible that $\alpha$ appears between necklaces other than $\beta_1$ and $\beta_2$ in $\bN_k(n,w^\uparrow)$.
\end{exam}

The following proof demonstrates that there is an efficient way to determine which necklaces contain $\alpha$ in $\bN_k(n,w^\uparrow)$. We put forward that the last necklace lexicographically smaller than $\beta_2$ in $\bN_k(n,w^\uparrow)$ must contain the suffix $\bp$, and the first necklace larger than $\beta_1$ in $\bN_k(n,w^\uparrow)$ must contain the prefix $\bq$. In other words, $\alpha = \bp \bq$ is found in the necklaces that encompass $\beta_1 \beta_2$ in $\bN_k(n,w^\uparrow)$.

\begin{proof} 
%
We can divide this into a proof by cases by considering the location of the increasing index between $\beta_1$ and $\beta_2$. There are two cases that must be considered when searching for $\alpha$. \\

\noindent
\textbf{Case 1.} Suppose that the increasing index between $\beta_1$ and $\beta_2$ is within the suffix $\bp$. Let us consider the positions of $\bp$ and $\bq$ in $\bN_k(n,w^\uparrow)$:

\begin{enumerate}
    \item $\bp$: By Lemma \ref{lem:lexorder}, $\beta_1$ and $\beta_2$ must share a common prefix up until the increasing index, both necklaces must share the common prefix of $\bq$. This means that $\beta_1 = \bq \bp$, with weight at least $w$ by definition of $\alpha$. Furthermore, we know that $\bp$ is not a maximal suffix, as it contains an increasing index. So, $\beta_1$ contains the suffix $\bp$ and must exist in the sequence $\bN_k(n,w^\uparrow)$.

    \item $\bq$: Let $g$ be the length of $\bq$ and let $\beta_3$ be a necklace of the form $\bq \cdot \tt{k}^{n-g}$. By Lemma \ref{lem:max}, $\beta_3$ must come after $\beta_1$ in lexicographic order.

    Clearly, $\beta_3$ must be heavier than $\beta_1$, thus $\beta_3$ is also a necklace in $\bN_k(n,w^\uparrow)$. Lexicographically, as $\beta_1$ and $\beta_3$ have the same prefix, $\bq$, no necklace between the two of them in $\bN_k(n)$ can have a differing prefix. If any necklace from $\bN_k(n)$ between $\beta_1$ and $\beta_3$ is included in $\bN_k(n,w^\uparrow)$, it also has the prefix $\bq$. Since $\beta_3$ is in $\bN_k(n,w^\uparrow)$, this means the the first necklace in $\bN_k(n,w^\uparrow)$ lexicographically larger than $\beta_1$ must have the prefix $\bq$.
\end{enumerate}

Therefore, in this case we can derive the location of $\alpha$ in $\bN_k(n,w^\uparrow)$ using $\bN_k(n)$. If $\alpha = \bp \bq$ in $\bN_k(n)$, where $\bp$ is the suffix of some necklace $\beta_1$ and $\bq$ is the prefix of some necklace $\beta_2$, then in $\bN_k(n,w^\uparrow)$ $\beta_1$ also appears in $\bN_k(n,w^\uparrow)$, and $\bq$ is found in the subsequent necklace. \\

\noindent
\textbf{Case 2.} Suppose that the increasing index between $\beta_1$ and $\beta_2$ is within the prefix $\bq$. Let us consider the positions of $\bp$ and $\bq$ in this case as well:

\begin{enumerate}
    \item $\bq$: By Lemma \ref{lem:lexorder}, $\beta_1$ must have a prefix lexicographically smaller than $\bq$. Additionally, Corollary \ref{cor:max2} tells us that $\bp$ must be a maximal suffix of the form $\tt{k}^g$. Similarly to the previous case, we can use Lemma \ref{lem:max} to construct a necklace $\beta_3$ which contains the prefix $\bq$ and the maximal suffix $\bp$.
    
    Let $\beta_3 = \bq \bp$. As $\alpha=\bp \bq$ has a weight of at least $w$, $\beta_3$ must also have a weight of at least $w$.Once again, we can observe that any necklace between $\beta_2$ and $\beta_3$ in $\bN_k(n)$ must contain the prefix $\bq$. Therefore, as $\beta_3$ is in $\bN_k(n,w^\uparrow)$, the necklace following $\beta_1$ in $\bN_k(n,w^\uparrow)$ has the prefix $\bq$.

    \item $\bp$: Unlike the first case, it is not certain that $\beta_1$ has a weight of $w$ or greater and may not be present in $\bN_k(n,w^\uparrow)$. However, suppose that there is some necklace in $\bN_k(n,w^\uparrow)$ smaller than $\beta_1$ lexicographically, but does not contain the suffix $\bp$. Since $\bp$ is a maximal suffix, we can construct another necklace, $\beta_0$, which has the same prefix as this necklace, but contains the maximal suffix $\bp$. 

    Let $\beta_0$ be a necklace that has a prefix lexicographically smaller than $\bq$, and containing the maximal suffix $\bp$. By Lemma \ref{lem:max}, this necklace exists. Note that $\beta_0$ is both heavier than the supposed previous necklace, and comes between that necklace and $\beta_2$ lexicographically, as $\beta_2$ necessarily contains a larger prefix. In other words, if any necklace exists before $\beta_1$ in $\bN_k(n,w^\uparrow)$ that does not contain the suffix $\bp$, there is always another necklace $\beta_0$ that contains $\bp$ as a suffix that comes before $\beta_2$. 

    \begin{exam}
    Suppose $n=4$, $k=3$, and $w=7$. Let $\alpha = \tt{\textcolor{blue}{33}\textcolor{red}{11}}$ $\beta_1 = \tt{0322}$, and $\beta_2 = \tt{\textcolor{red}{11}23}$. $\beta_1$ comes before $\beta_2$ in $\bN_k(n,w^\uparrow)$. As $\tt{\textcolor{blue}{33}}$ is a max suffix, we can construct the following necklace with the same prefix as $\beta_1$:
    %
    \[ \beta_0 = \tt{03\textcolor{blue}{33}} \]
    %
    This necklace appears between $\beta_1$ and $\beta_2$, and contains the maximal suffix found in $\alpha$. This is true regardless of the prefix of $\beta_1$. In this case, the sequence becomes:
    
    \[ \tt{\cdots, 0322, \ldots, 03\textcolor{blue}{33}, \ldots, \textcolor{red}{11}23}, \ldots \]
    \end{exam}
    
    If no necklace exists in $\bN_k(n,w^\uparrow)$ before $\beta_2$, then recall Property \ref{prop:last}. When considered cyclically the previous necklaces in $\bN_k(n,w^\uparrow)$ are $(k-1) \tt{k}^{n-1}$ and $k^n$. Since the length of $\bp$, $g$, is less than $n$, the previous suffix would once again be $\bp = \tt{k}^g$. In all cases, the suffix preceding $\beta_2$ is $\bp$.

\end{enumerate}

% Step three: put it together

Therefore, in both cases we derive the position of $\alpha = \bp \bq$ in $\bN_k(n,w^\uparrow)$ using the position of $\alpha$ in $\bN_k(n)$. Regardless of the increasing index of $\alpha$, the prefix $\bq$ is always contained in the first necklace in $\bN_k(n)$ after $\beta_1$ that has a weight of at least $w$.
%
\end{proof}

The algorithm outlined in Section \ref{sec:granddaddy} locates $\beta_2$ in $\mathcal{O}(n^2)$ time using $\mathcal{O}(n)$ space. The previous proof shows us $\bq$ is contained in the first necklace in $\bN_k(n,w^\uparrow)$ after $\beta_1$ lexicographically. The following algorithm demonstrates how this is possible in $\mathcal{O}(n)$ time using $\mathcal{O}(n)$ space.

\begin{lemma}
    Let $\beta_1 = \tt{a}_1 \tt{a}_2 \cdots \tt{a}_n$ be a necklace in $\bN_k(n)$ that contains the prefix $\bq$. The first necklace $\beta_2$ with a weight of at least $w$ following $\beta_1$ in $\bN_k(n)$ takes the form $\beta_2 = \tt{a}_1 \tt{a}_2 \cdots \tt{a}_{i-1} \cdot \tt{x} \cdot \tt{k}^{n-i}$, where $i$ is the largest possible value such that the weight of $\beta_2$ equals $w$.
\end{lemma}

\begin{proof}
    Clearly, the necklace in $\bN_k(n,w^\uparrow)$ closest after $\beta_1$ lexicographically, $\beta_2$, will have an increasing index in the highest possible position. If a necklace of the required weight shared a larger common prefix with $\beta_1$, then that necklace would come before $\beta_2$. Likewise, the value $x$ at the increasing index $i$ must have the smallest possible value so that the weight of $\beta_2$ is at least $w$. Recall from Lemma \ref{lem:max} that increasing the last non-$k$ character always results in a valid necklace. Thus, the next necklace lexicographically can be constructed by repeatedly incrementing that final character. Notice that this will always maintain the longest possible common prefix between $\beta_1$ and $\beta_2$ by the time the weight $w$ is reached.
\end{proof}

By this lemma, the algorithm for finding the next valid necklace $\beta_2$ in $\bN_k(n,w^\uparrow)$ given a necklace $\beta_1$ is simple. Given $\beta_1$ and the required weight $w$, we increment the final non-maximal character in the sequence. If incrementing this character does not meet the weight requirement, then we repeat this until the weight of the sequence reaches $w$. Each character can be increased in linear time, and at most $n$ characters need to be modified. Thus, this algorithm works in $\mathcal{O}(n)$ time and $\mathcal{O}(n)$ space.

Therefore, given some string  $\alpha = \bp\bq$ in $\bS_k(n)$ with weight at least $w$ such that $\bq$ is the longest suffix of $\alpha$ such that $\bq\bp$ is a necklace, we can find the necklace $\beta$ in $\bN_k(n,w^\uparrow)$ that contains $\bq$ in $\mathcal{O}(n^2)$ time using $\mathcal{O}(n)$ space. We do this by applying the algorithm outlined in~\cite{kociumaka,ranking} to find $\beta$ in $\bN_k(n)$, then find the first necklace lexicographically larger than $\beta$ with a weight of at least $w$.






\end{comment}

\end{document}
